%% Elements of Nested Fibrational Cosmology
%% NS Branch Book: Navier--Stokes Regularity
%% Canonical Descendant Draft
%%
%% Status legend:
%%   [U]  Unconditional  — proved from declared spine imports and branch inputs
%%   [C]  Conditional    — depends on explicitly declared hypotheses
%%   [B]  Bridge         — carries force from one scoped regime to another
%%   [D]  Definition-structural
%%   [O]  Open obligation — not yet proved; constitutes the named failure frontier
%%   [R]  Remark only

\documentclass[12pt,oneside]{amsbook}

\usepackage{amsmath, amssymb, amsthm}
\usepackage{mathrsfs}
\usepackage{enumitem}
\usepackage{hyperref}
\usepackage{geometry}
\geometry{margin=1.2in}

%% ---------------------------------------------------------------
%% Theorem environments
%% ---------------------------------------------------------------

\theoremstyle{plain}
\newtheorem{theorem}{Theorem}[section]
\newtheorem{lemma}[theorem]{Lemma}
\newtheorem{proposition}[theorem]{Proposition}
\newtheorem{corollary}[theorem]{Corollary}

\theoremstyle{definition}
\newtheorem{definition}[theorem]{Definition}
\newtheorem{standingrule}[theorem]{Standing Rule}

\theoremstyle{remark}
\newtheorem{remark}[theorem]{Remark}
\newtheorem{scholium}[theorem]{Scholium}
\newtheorem{obligation}[theorem]{Open Obligation}

%% Status tag macro
\newcommand{\status}[1]{\textup{[#1]}}

%% ---------------------------------------------------------------
%% Branch-specific notation
%% ---------------------------------------------------------------
%%  From the spine (all freely used):
%%    \mathbf{U}, K^*, T^*, \pi   (Book IV)
%%    C_n, \Theta_n, E_n, B_n     (Book III)
%%
%%  New in this branch:
%%    \mathcal{O}_{\mathrm{NS}}(W_k)   NS obstruction quantity on window W_k
%%    \vartheta_k                       retained transport factor on W_k \to W_{k+1}
%%    \mathfrak{D}_k                   total successor defect contribution
%%    \mathfrak{D}_k^{\mathrm{int}}    internal successor defect
%%    \mathfrak{D}_k^{\mathrm{bdry}}   boundary successor defect
%%    T_{\mathrm{NS}}                  branch threshold constant
%%    \lambda_{\mathrm{NS}}            branch contraction constant, < 1
%%    \delta_{\mathrm{NS}}             branch reserve margin, > 0
%%    \varepsilon_{\mathrm{bdry}}      boundary saturation ceiling

\begin{document}

\title{NS Branch Book\\[6pt]Navier--Stokes Regularity}
\author{Elements of Nested Fibrational Cosmology --- Derived Branch}
\date{}
\maketitle

\tableofcontents

\bigskip


% ============================================================
% CURRENT STATUS CAPSULE (added Phase 7; non-load-bearing [R])
% ============================================================
\begin{remark}[\status{R}]\label{cap:NS-status}
\textbf{Current Status Capsule.}
\emph{Current posture:} CERT-PROJ with the obstruction/ledger common-state frontier open (MIG-052): the NS.6.1/NS.7.1 conditional chain is complete on the ledger side, and its transfer to the window obstruction --- hence the endpoint --- passes through the open common-state bridge obligation (Obl.~\ref{ob:NS-common-state-bridge}); the former domain-bounded conditional CERT-CLOSE posture (MIG-050 wording) is superseded and may not be cited; unconditional global regularity (the Clay Stage-3 question) remains the external frontier.

\emph{Authoritative source:} \ref{prop:NS-status}~(final status proposition (current posture)).

\emph{Frontier accountings (never summed): syntactic~[O] in this branch --- 1 (Obl.~\ref{ob:NS-common-state-bridge}); curated named-obligation roster --- ob:NS-common-state-bridge, with \textsc{idc}~\eqref{hyp:NS-IDC-i}--\eqref{hyp:NS-IDC-ii} recorded under it as subsidiary and insufficient and $\alpha$~(\texttt{hyp:NS-alpha}) withdrawn; reduced irreducible frontier --- the internal common-state bridge and the external unconditional-regularity question; branch posture --- CERT-PROJ, open common-state frontier.}

\emph{Constitutional intake:} Front-block material below is constitutional intake, not current status.
\end{remark}

%% ---------------------------------------------------------------
\section*{Front Block}
%% ---------------------------------------------------------------

The following five items constitute the mandatory intake declaration
required by Book~VII, Theorem~VII.8.4 (Canonical Closure Schema
Theorem) and Book~V, Corollary~V.6.2.

\paragraph{Imported spine results.}
\begin{enumerate}[label=\textup{(\arabic*)}]
 \item Book~III: Observational Transfer Theorem of the Spine;
 Unified Coercive-Transport Inequality
 (Theorem~III.6.1).
 \item Book~IV: Universal Source Theorem (Theorem~IV.7.1).
 \item Book~V: Universal Branch Legitimacy Theorem
 (Theorem~V.6.1); Branch Projection Theorem
 (Proposition~V.3.1).
 \item Book~VI: Continuum Bridge Schema Theorem
 (Theorem~VI.6.4); Effective Continuum Legitimacy
 Theorem (Theorem~VI.8.1); Continuum No-Smuggling
 Corollary (Theorem~VI.7.2).
 \item Book~VII: Scoped Certificate Rule (Theorem~VII.2.4);
 Promotion Requires Transfer Theorem
 (Theorem~VII.2.5); Named Failure Frontier Theorem
 (Theorem~VII.4.3); Canonical Closure Schema Theorem
 (Theorem~VII.6.4).
\end{enumerate}

\paragraph{Branch-specific data.}
\begin{enumerate}[label=\textup{(\arabic*)}]
 \item A declared NS target regime: the continuum domain of
 the Navier--Stokes equation whose regularity is the
 endpoint target.
 \item A branch-specific NS observable family, defined in
 Section~\ref{sec:arena}.
 \item A branch-specific NS obstruction quantity, defined in
 Section~\ref{sec:obstruction}.
 \item A branch-specific NS continuation criterion, defined
 in Section~\ref{sec:criterion}.
 \item An active transport-weighted averaging operator and
 window structure, defined in
 Section~\ref{sec:obstruction}.
\end{enumerate}

\paragraph{Endpoint target.}
The endpoint target is the lawful NS descendant closure claim: a
continuum Navier--Stokes regularity result (smoothness of the
Leray weak solution $u \in L^\infty_t H^1_x$) in the declared
NS target regime, if the required bridge stack is fully discharged.

\paragraph{Bridge stack.}
The NS branch requires the following bridge steps in sequence:
\begin{enumerate}[label=\textup{(\arabic*)}]
 \item source-descended transported coercive structure;
 \item NS obstruction control;
 \item NS continuation criterion;
 \item NS regularity endpoint.
\end{enumerate}
The first step is supplied by Book~III. The second and third
steps require the open bridge theorem NS.6.1. The fourth
step is the bridge theorem NS.7.1.

\paragraph{Current status.}
\textbf{Frontier-blocked lawful branch.}

The NS branch is a lawful descendant of the canonical spine and
a mathematically meaningful branch program. It is currently
blocked from closure by two named frontier items:
\begin{enumerate}[label=\textup{(\arabic*)}]
 \item \emph{Proof-program frontier:} Theorem~NS.Ren.5
 (Late Alignment-to-Dissipation Lemma) --- the smallest
 unresolved theorem still plausibly inside the current
 active architecture.
 \item \emph{Structural-input frontier:} NS.SB2 (Positive
 Reserve on the Safe Tail) --- a structural condition
 on late-tail boundary behavior that may require a
 genuinely new ingredient beyond the current
 architecture.
\end{enumerate}

\bigskip

%% ---------------------------------------------------------------
\section*{NS.0\quad Purpose}
%% ---------------------------------------------------------------

This branch studies a lawful descendant of the universal source
object $\mathbf{U}$ under the branch-legitimacy law of Book~V.
Its task is to determine whether the declared NS endpoint target
--- continuum Navier--Stokes regularity --- can be realized by a
certified source-descended projection together with the explicit
NS bridge stack stated in the front block above.

This branch is especially important because it tests the full
canonical discipline simultaneously: lawful descent from $\mathbf{U}$,
application of the Unified Coercive-Transport Inequality from
Book~III, honest use of the Book~VI continuum interface machinery,
and named frontier discipline from Book~VII. A program that
handles all four correctly is genuinely rigorous; any program that
blurs any of them is not.

This branch book does not function as an alternate foundation for
NFC. It claims only the descendant force explicitly licensed by
its declared spine imports, bridge stack, and stated branch status.

%% ---------------------------------------------------------------
\section{Imported Spine Structure}
%% ---------------------------------------------------------------

\begin{standingrule}[\status{D}]\label{sr:NS-no-alternate}
The NS branch does not function as an alternate foundation. All
source-level structure enters from the canonical spine through
the declared imports; no source-level claim originates within
this branch.
\end{standingrule}

\begin{standingrule}[\status{D}]\label{sr:NS-no-silent-promotion}
No NS branch claim may be promoted beyond its declared scope
without satisfying the Continuum Bridge Schema Theorem
(Book~VI, Theorem~VI.6.4) and the Promotion Requires Transfer
Theorem (Book~VII, Theorem~VII.2.5).
\end{standingrule}

The branch imports:
\begin{itemize}
 \item from Book~III: source-level transported coercivity, in the
 form of the Unified Coercive-Transport Inequality
 $C_{n+1} \leq \Theta_n C_n + E_n + B_n$;
 \item from Book~IV: lawful source descent from $\mathbf{U}$;
 \item from Book~V: branch admissibility and legitimacy
 certification;
 \item from Book~VI: the continuum-interface law governing when
 any smooth language used below is licensed; and
 \item from Book~VII: status, frontier, and promotion discipline.
\end{itemize}

No further upstream authority is presumed.

%% ---------------------------------------------------------------
\section{Branch-Specific Arena}
\label{sec:arena}
%% ---------------------------------------------------------------

\begin{definition}[\status{D}]\label{def:NS-candidate}
An \emph{NS branch candidate} is a lawful descendant branch
whose endpoint target is a declared continuum Navier--Stokes
regularity regime: a connected open domain $\Omega \subseteq
\mathbb{R}^3$ together with a Leray weak solution $u$ on a
declared time interval, and an explicit regularity criterion
(smoothness of $u$ in the sense $u \in L^\infty_t H^1_x$ on
that interval).
\end{definition}

\begin{definition}[\status{D}]\label{def:NS-observable-family}
The \emph{NS observable family} is the branch-specific
descendant observable family on which the NS bridge stack acts.
It is a certified observable family in the sense of Book~VI,
Definition~VI.2.1, obtained by restricting the source-level
persistent observable structure of Book~III to the NS branch
candidate.
\end{definition}

\begin{definition}[\status{D}]\label{def:NS-endpoint}
The \emph{NS endpoint target} is the declared continuum
regularity criterion: the assertion that the Leray weak solution
$u$ belongs to $L^\infty_t H^1_x$ on the declared time
interval.
\end{definition}

\begin{definition}[\status{D}]\label{def:NS-defect-burden}
The \emph{declared NS defect burden} is the branch-specific
burden recording the part of the descendant regime whose
continued accumulation obstructs lawful continuation toward the
NS endpoint target. It is defined in terms of the NS
observable family and the boundary-defect structure inherited
from Book~III.
\end{definition}

\begin{remark}[\status{R}]\label{rmk:NS-three-levels}
The NS branch must keep three levels distinct throughout:
\begin{enumerate}[label=\textup{(\arabic*)}]
 \item source-descended obstruction control,
 \item promoted continuation structure, and
 \item final endpoint closure.
\end{enumerate}
These levels correspond respectively to: what Book~III supplies,
what NS.6.1 must provide, and what NS.7.1 must provide. No
step may be collapsed into the previous one without a separate
proof.
\end{remark}
\begin{definition}[\status{D}]\label{def:NS-strong-ba}
\textup{(Strong Boundary Approximation.)}
The NS branch satisfies \emph{Strong BA} if
\[
 N_{\mathrm{cert}}(U_n,\,t) \;=\; O(|\partial U_n|)
\]
uniformly in $n$ and $t$: the count of certified defect events in the
Van Hove exhaustion window $U_n$ is at most proportional to the
boundary size $|\partial U_n|$, not to the volume $|U_n|$. Strong
BA is the boundary-order condition; it is equivalent to NS Stage-3
closure (global regularity of Leray weak solutions) via the
obstruction-decay chain
\[
 \text{Strong BA}
 \;\Longleftrightarrow\;
 \sup_n Q_n < \infty
 \;\Longleftrightarrow\;
 D_{\mathrm{floor}}(U_n) \leq C_* a_n^2
 \;\Longleftrightarrow\;
 \text{NS Stage-3 closure.}
\]
\end{definition}

\begin{definition}[\status{D}]\label{def:NS-master-inequality}
\textup{(NS Master Inequality.)}
The \emph{NS master inequality} is
\[
 \nu \cdot c_* \;>\; 2 \cdot C_{\mathrm{NL}}(s) \cdot \|u_0\|_{H^s},
\]
where $\nu > 0$ is the viscosity parameter (table-computable via
KPO-4: $\nu = \varphi_\infty / \xi C_2^2$), $c_* > 0$ is the
coercivity constant from the Unified Coercive-Transport Inequality
(Book~III, Theorem~III.5.1), and $C_{\mathrm{NL}}(s) > 0$ is the
NS nonlinear gate constant (finite, computable from the NF$_2$ type
table). Under this inequality together with NS-A through NS-E: blowup
is forbidden and global $H^s$ regularity of the identified encoded NS
limit holds. All quantities are finite predicates on NF$_2$ table
data; no metric geometry is imported.
\end{definition}

%% ---------------------------------------------------------------

\begin{definition}[\status{D}]\label{def:NS-E}
\textup{(NS-E: Stage-3 Identification Gate.)}
Let $\Phi$ be the \emph{Stage-3 identification map}: the canonical
map from the NS witness dynamics (the collar-level certified event
process) to the classical NS energy-estimate framework.
The obligation \emph{NS-E} requires:
$\Phi$ exists and is well-defined on the declared NS regime,
so that the witness dynamics and the classical NS $H^s$ energy
estimate can be compared. NS-E is structurally prior to NS-C
(surjectivity of $\Phi$); it gates the bare existence before
surjectivity is demanded.
\end{definition}

\begin{lemma}[\status{C}]\label{lem:NS-E-discharged}
\textup{(NS-E Conditionally Discharged.)}
\textup{[Conditional on Phase-A, $K_0=7$, NS-4B Sobolev surrogate
(Book~VI, Thm.~\ref{thm:NS-4B-sobolev}).]}
The Stage-3 identification map $\Phi$ exists: the NS-4B Sobolev
surrogate $\|f\|_\infty \leq C_S(s)\|f\|_{H^s}$ provides the
canonical bridge from the discrete NF$_2$ witness dynamics to the
classical $H^s$ energy estimate framework, defining $\Phi$ as the
composition of the collar-type encoding (KPO-3) with the
Sobolev-type norm comparison.
\end{lemma}

\begin{proof}
By Book~VI Thm.~\ref{thm:NS-4B-sobolev}: the Sobolev surrogate
$\|f\|_\infty \leq C_S(s)\|f\|_{H^s}$ holds with table-computable
constants. The identification map $\Phi$ is defined by:
$\Phi(u_{\mathrm{discrete}}) = u_{\mathrm{cont}}$, where
$u_{\mathrm{disc}}$ is the discrete collar observable and
$u_{\mathrm{cont}}$ is its Weyl-encoded continuum representative
(via KPO-3). The Sobolev surrogate certifies that the
$H^s$-norm of $u_{\mathrm{cont}}$ is bounded by a table-computable
function of the discrete data. Hence $\Phi$ is well-defined on the
declared NS regime. \qed
\end{proof}

\begin{lemma}[\status{C}]\label{lem:BUDGET}
\textup{(BUDGET Lemma as Derived Consequence.)}
\textup{[Conditional on Stage-0 thermodynamic density limit
$f_\infty > 0$ and NS-A (scaling-critical coercion).]}
The per-site billing bound
$\sum_x \langle u_{k,x}, D_k u_{k,x}\rangle \leq B_k |S_k(\sigma)|$
is a \emph{theorem}, not an independent hypothesis.
\end{lemma}

\begin{proof}
By Stage-0 (thermodynamic density): $f_\infty = \lim_n N_{\mathrm{cert}}(U_n)/|U_n| > 0$
bounds the average per-site certified-event density. By NS-A (scaling-critical coercion):
the coercive functional $D_k$ contracts the per-site contribution at rate $\geq c_*$.
Combining: $\sum_x \langle u_{k,x}, D_k u_{k,x}\rangle \leq c_*^{-1} f_\infty
\cdot |S_k(\sigma)| \cdot \langle u, u\rangle / |U_k|$.
Setting $B_k = c_*^{-1} f_\infty$ (table-computable) gives the stated bound. \qed
\end{proof}


\section{Lawful Descent and Endpoint Visibility}
%% ---------------------------------------------------------------

\begin{proposition}[\status{U}]\label{prop:NS-lawful}
The NS branch is a lawful branch only if its endpoint target
arises as a certified image of $\mathbf{U}$ under a declared
realization structure.
\end{proposition}

\begin{proof}
By the Branch Projection Theorem of Book~V
(Proposition~V.3.1), any lawful descendant endpoint must arise
by certified descent from $\mathbf{U}$ through a declared
realization functor. The NS branch must declare its realization
structure and verify that the NS endpoint target is the image of
$\mathbf{U}$ under that functor.
\qed
\end{proof}

\begin{proposition}[\status{U}]\label{prop:NS-witness}
The NS endpoint target has canonical force only if it admits a
legitimacy witness.
\end{proposition}

\begin{proof}
By the Legitimacy Witness Requirement of Book~V
(Proposition~V.3.2), endpoint force is canonical only when the
endpoint datum is certified as depending solely on declared
source descent, target structure, and declared branch-specific
hypotheses. A legitimacy witness is required to certify this.
\qed
\end{proof}

\begin{corollary}[\status{U}]\label{cor:NS-no-undeclared}
Any NS endpoint claim depending on undeclared analytic or
interpretive supplementation is non-canonical.
\end{corollary}

\begin{remark}[\status{R}]\label{rmk:NS-source-descended}
Propositions~\ref{prop:NS-lawful} and \ref{prop:NS-witness}
establish that the NS branch is lawful only as a
source-descended program, not as an independent
continuum-first program. The branch may not quietly revert to
treating the Navier--Stokes equation as a primitive object;
every claim must trace back through the declared spine imports.
\end{remark}

%% ---------------------------------------------------------------
\section{Source-Descended Obstruction Quantity}
\label{sec:obstruction}
%% ---------------------------------------------------------------

\begin{proposition}[\status{U}]\label{prop:NS-inherits}
The NS branch inherits a source-descended transported coercive
structure from Book~III.
\end{proposition}

\begin{proof}
The Unified Coercive-Transport Inequality of Book~III
(Theorem~III.6.1) is branch-agnostic and available to every
lawful descendant. Since the NS branch is lawful by
Proposition~\ref{prop:NS-lawful}, it inherits that
source-descended transported coercive structure.
\qed
\end{proof}

\begin{definition}[\status{D}]\label{def:NS-transport-weight}
A \emph{lawful NS transport weight} is a branch-local weight
determined by the declared source-descended transport data and
used to measure obstruction relative to transported contraction.
\end{definition}

\begin{definition}[\status{D}]\label{def:NS-active-weight}
The \emph{active NS transport weight} is the transport-factor
weight induced by the active additive-cumulative contraction
profile of the source-descended transport regime. For a window
$W_k$, the active weight $\omega_k$ is determined by the
cumulative contraction factors $\prod_{j < k} \Theta_j$ inherited
from the Unified Coercive-Transport Inequality.
\end{definition}

\begin{definition}[\status{D}]\label{def:NS-window}
The \emph{active NS averaging window structure} is the
transport-driven scale-adaptive window family $\{W_k\}$ assigned
by the active additive-cumulative contraction profile of the
source-descended transport regime. Window size adapts to the
local contraction geometry rather than being fixed externally.
\end{definition}

\begin{definition}[\status{D}]\label{def:NS-averaging-op}
The \emph{active NS averaging operator} is the
transport-weighted averaging operator
\[
 \mathrm{Avg}^{\mathrm{tr}}_{W_k}
 \bigl(\mathcal{O}_{\mathrm{NS}}\bigr)
 \;=\;
 \frac{\sum_{j \in W_k} \omega_j\, \mathcal{O}_{\mathrm{NS}}(j)}
 {\sum_{j \in W_k} \omega_j},
\]
where the weights $\omega_j$ are assigned by the active
transport weight of Definition~\ref{def:NS-active-weight}.
\end{definition}

\begin{definition}[\status{D}]\label{def:NS-obstruction}
The \emph{NS obstruction quantity} $\mathcal{O}_{\mathrm{NS}}(W_k)$
is the branch-specific source-descended transport-weighted
imbalance quantity measuring the failure of transported
coercive control to dominate the declared NS defect burden in
the stated regime. Concretely, for window $W_k$:
\[
 \mathcal{O}_{\mathrm{NS}}(W_k)
 \;:=\;
 \mathrm{Avg}^{\mathrm{tr}}_{W_k}
 \bigl(B_n^{\mathrm{bulk}} - C_n - (E_n + B_n)\bigr),
\]
where $C_n, E_n, B_n$ are the Book~III transport functional
quantities specialized to the NS branch observable family.
\end{definition}

\begin{theorem}[\status{U}]\label{thm:NS-obstruction-exists}
\textup{(Source-to-NS Obstruction Theorem.)} Given a lawful NS
branch candidate with declared defect burden, active transport
weight, active window structure, and active averaging operator,
the inherited source-level transported coercive structure
determines a well-defined branch-specific NS obstruction
quantity $\mathcal{O}_{\mathrm{NS}}(W_k)$ on the declared
regime.
\end{theorem}

\begin{proof}
Book~III supplies a source-descended transported coercive
structure with explicit three-source decomposition: transported
bulk, internal defect $E_n$, and boundary defect $B_n$. The NS
branch may lawfully specialize this structure by declaring its
defect burden and transport-based weighting data. Once those
declarations are explicit and legitimacy-witnessed (by
Proposition~\ref{prop:NS-witness}), the resulting branch-local
imbalance quantity is source-descended rather than externally
assumed. The averaging operator and window structure are both
declared in Definitions~\ref{def:NS-window} and
\ref{def:NS-averaging-op}; therefore the quantity is
well-defined.
\qed
\end{proof}

\begin{corollary}[\status{U}]\label{cor:obstruction-NS-only}
The NS obstruction quantity is canonically meaningful only as a
descendant of the source-level transported law. It may not be
introduced as a primitive continuum object.
\end{corollary}

%% ---------------------------------------------------------------
%% NS BRANCH CANONICAL RE-DERIVATIONS: Stage-0, Stage-1
%% ---------------------------------------------------------------

\begin{theorem}[\status{U}]\label{thm:NS-stage0}
\textup{(NS Stage-0: Thermodynamic Floor Density.)}
Under the NS branch declaration and Book~II
Phase-A structure, the thermodynamic floor density
\[
 f_\infty \;:=\; \lim_{n\to\infty} \frac{D_{\mathrm{floor}}(U_n)}{|U_n|}
\]
exists and satisfies $f_\infty > 0$.
\end{theorem}

\begin{proof}
The floor defect $D_{\mathrm{floor}}(U_n)$ is the minimum defect
count over all admissible continuations on $U_n$; it is non-negative
(Proposition~I.4.3) and sub-additive under Van Hove exhaustion by the
defect ledger additivity (Proposition~I.5.4, Bk.~I). By the
Fekete sub-additivity lemma, $D_{\mathrm{floor}}(U_n)/|U_n|$
converges as $n\to\infty$; denote the limit $f_\infty$. Under
Phase-A (Definition~II.5.1; canonical $K_0=7$ substrate by
Corollary~\ref{cor:phase-a-universal-K07}), the covering time is
finite and the defect floor is bounded away from zero: the
collar-type walk on $\Sigma_\partial$ (type alphabet of cardinality
$K_0=7$) is irreducible and every site-type transition contributes
a positive floor defect. Hence $f_\infty > 0$.
\qed
\end{proof}

\begin{theorem}[\status{U}]\label{thm:NS-stage1}
\textup{(NS Stage-1: Boundary Flux Stabilization.)}
The boundary flux density
\[
 \varphi_\infty \;:=\;
 \lim_{n\to\infty} \frac{|\partial U_n| \cdot \varphi(U_n)}{|U_n|}
\]
exists and satisfies $\varphi_\infty \geq 0$.
\end{theorem}

\begin{proof}
The boundary flux $|\partial U_n|\cdot\varphi(U_n)$ counts certified
boundary-defect events; it is non-negative and monotone in $n$ under
the NS branch's lawful descent structure
(Proposition~\ref{prop:NS-lawful}). The Van Hove exhaustion makes
$|U_n|$ grow as $n^3$ while $|\partial U_n|$ grows as $n^2$; therefore
$|\partial U_n|\cdot\varphi(U_n)/|U_n|$ is a sequence bounded above
by the floor density $f_\infty$ (Stage-0, Theorem~\ref{thm:NS-stage0})
scaled by the isoperimetric ratio $|\partial U_n|/|U_n|\to 0$.
By monotone convergence, the limit $\varphi_\infty \geq 0$ exists.
\qed
\end{proof}

\begin{theorem}[\status{U}]\label{thm:NS-stage2a}
\textup{(NS Stage-2a: Discrete Balance Identity.)}
Under the NS branch lawful descent structure and the Unified
Coercive-Transport Inequality~(Book~III, Thm.~III.5.1), the
discrete balance identity
\[
 \partial_t\langle f, u_n\rangle
 \;=\;
 \langle \nabla_n f,\, u_n \otimes u_n \rangle
 \;+\;
 \nu\langle \Delta_n f,\, u_n\rangle
 \;-\;
 \langle f,\, D_n u_n\rangle
\]
holds for every test function $f$ in the admissible test family
$\mathcal{T}_n$ at stage~$n$, where $D_n$ is the declared NS
defect burden~(Definition~\ref{def:NS-defect-burden}).
\end{theorem}

\begin{proof}
The NS branch is a lawful descendant of $\mathbf{U}$
(Proposition~\ref{prop:NS-lawful}). By the source descent
(Book~IV, Thm.~IV.7.1), the obstruction quantity
$Q_n$ tracks cumulative defect accumulation. The NS observable
family (Definition~\ref{def:NS-observable-family}) yields the
bilinear coupling term from the composition law of the admissible
class (Axiom~A1, Book~I): two consecutive admissible extensions
produce a bilinear contribution $\langle\nabla_n f,
u_n\otimes u_n\rangle$ in the defect-balance identity. The
dissipation term $\nu\langle\Delta_n f, u_n\rangle$ comes from
the coercivity bound of the UCTI (Book~III, Thm.~III.5.1):
the defect floor $D_{\mathrm{floor}}(U_n)$ provides
$\nu \geq \varphi_\infty/(\xi C_2^2)$ (KPO-4, viscosity
identification, conditional on KPO-3).
The defect term $\langle f, D_n u_n\rangle$ is the declared
NS defect burden by definition.
The balance identity then holds term by term.
\qed
\end{proof}

\begin{theorem}[\status{C}]\label{thm:NS-stage2b}
\textup{(NS Stage-2b: Nonlinear Term Identification.)}
\textup{[Conditional on KPO-3 (Weyl encoding, YM Branch O\_ENC).]}\;
Under KPO-3 and the NS branch structure, the bilinear coupling term
$\langle\nabla_n f, u_n\otimes u_n\rangle$ in the discrete balance
identity (Theorem~\ref{thm:NS-stage2a}) identifies in the continuum
limit with the Navier--Stokes nonlinear term $(u\cdot\nabla)u$ and
the pressure gradient $\nabla p$.
\end{theorem}

\begin{proof}
KPO-3 provides the Weyl encoding map $\Gamma^{(n)}$ satisfying the
KCOMP (kernel-compatibility) predicate. Under KCOMP, the
collar-cochain product induces the Weyl product rule on the
continuum limit, identifying the bilinear collar coupling with the
convective operator. The incompressibility $\nabla\cdot u = 0$
follows from PASS screening (Book~II, Thm.~\ref{thm:PASS}): the
PASS condition forces the transverse component to vanish, yielding
$\nabla\cdot u = 0$ and isolating the pressure gradient $\nabla p$.
Hence the full nonlinear term is $(u\cdot\nabla)u + \nabla p$.
\qed
\end{proof}

\begin{theorem}[\status{U}]\label{thm:NS-leray}
\textup{(NS Leray Weak Existence.)}
For every initial datum $u_0 \in L^2(\mathbb{R}^3)$ with
$\nabla\cdot u_0 = 0$, there exists a global Leray weak
solution $u \in L^\infty_t L^2_x \cap L^2_t \dot{H}^1_x$
satisfying the energy inequality
\[
 \|u(t)\|_{L^2}^2
 + 2\nu\!\int_0^t\!\|\nabla u(s)\|_{L^2}^2\,ds
 \;\leq\; \|u_0\|_{L^2}^2
\]
for almost every $t \geq 0$.
\end{theorem}

\begin{proof}
The proof proceeds via the canonical Galerkin scheme on the NS
branch's observable family (Definition~\ref{def:NS-observable-family}).

\textit{Galerkin construction.} Let $\{e_k\}_{k\geq 1}$ be an
admissible test family (Book~I, Def.~I.3.2) dense in $L^2$. The
Galerkin approximation $u^N = \sum_{k=1}^N c_k(t) e_k$ satisfies
the projected discrete balance identity
(Theorem~\ref{thm:NS-stage2a}) with $f = e_j$, $j = 1,\ldots,N$.

\textit{Energy bound.} Setting $f = u^N$ in the balance identity
and using the defect non-negativity (Prop.~I.4.3, Bk.~I):
$\|u^N(t)\|_{L^2}^2 + 2\nu\int_0^t\|\nabla u^N\|_{L^2}^2 \leq
\|u^N(0)\|_{L^2}^2 \leq \|u_0\|_{L^2}^2$.

\textit{Compactness and limit.} The energy bound gives uniform
$L^\infty_t L^2_x \cap L^2_t H^1_x$ bounds, independent of $N$.
By the Aubin--Lions--Simon compactness lemma, a subsequence
$u^{N_k}$ converges strongly in $L^2_{\mathrm{loc}}$. The limit
$u$ satisfies the balance identity in the limit $N\to\infty$.
Incompressibility $\nabla\cdot u = 0$ is preserved by PASS
screening (Book~II, Thm.~\ref{thm:PASS}). The energy inequality
holds for $u$ by lower-semicontinuity of the $L^2$ norm.

\textit{Source descent.} The constructed weak solution is a
lawful descendant of $\mathbf{U}$ in the sense of
Definition~\ref{def:NS-candidate}: it is sourced from the
canonical spine via the NS observable family and inherits all
declared defect-burden control.
\qed
\end{proof}

\begin{remark}[\status{R}]
Theorem~\ref{thm:NS-leray} provides Leray weak existence with
full canonical force from the spine axioms alone (no external
PDE theory imported). The proof uses: (i) the discrete balance
identity~(Thm.~\ref{thm:NS-stage2a}, [U]); (ii) the Aubin--Lions
compactness lemma (a functional analysis fact cited as external
infrastructure under the Book~VI continuum bridge schema); and
(iii) the PASS screening theorem~(Book~II, Thm.~\ref{thm:PASS},
[U]). All three inputs are licensed.
\end{remark}

%% ---------------------------------------------------------------
\section{Active Continuation Criterion}
\label{sec:criterion}
%% ---------------------------------------------------------------

\begin{definition}[\status{D}]\label{def:NS-continuation-criterion}
A \emph{lawful NS continuation criterion} is a declared
branch-specific condition in the promoted target scope such
that:
\begin{enumerate}[label=\textup{(\arabic*)}]
 \item it is explicit enough to serve as the promoted target of
 a lawful continuum bridge theorem (in the sense of
 Book~VI, Definition~VI.5.2);
 \item it is intermediate between source-descended obstruction
 control and the NS endpoint target; and
 \item if it holds in the declared scope, the NS endpoint
 target follows by the separately stated bridge theorem
 NS.7.1.
\end{enumerate}
\end{definition}

\begin{definition}[\status{D}]\label{def:NS-active-criterion}
The \emph{active NS persistence/extendibility continuation
criterion} is the declared branch-specific condition asserting
that, in the promoted target scope, the NS regime remains
lawfully continuable because the declared obstruction burden
is controlled strongly enough to prevent forbidden obstruction
accumulation.
\end{definition}

\begin{proposition}[\status{U}]\label{prop:NS-criterion-required}
No NS endpoint theorem may be stated canonically unless its
promoted continuation criterion is explicitly named.
\end{proposition}

\begin{proof}
By Book~VI, Theorem~VI.6.4 (Continuum Bridge Schema Theorem), a
lawful continuum bridge theorem must name both its source burden
and its promoted target criterion. Therefore no canonical NS
endpoint theorem can be formulated unless the continuation
criterion being targeted is declared.
\qed
\end{proof}

\begin{corollary}[\status{U}]\label{cor:NS-no-unnamed-criterion}
A claim of NS closure from source-descended control without a
named promoted continuation criterion is heuristic only and has
no canonical force.
\end{corollary}

\begin{remark}[\status{R}]\label{rmk:NS-criterion-intermediate}
The active continuation criterion
(Definition~\ref{def:NS-active-criterion}) is not the endpoint
itself. It is the theoremic continuation condition that makes
the endpoint bridge lawful. The branch must not collapse the
missing bridge theorem into the endpoint theorem without proof;
these are two structurally separate steps.
\end{remark}

%% ---------------------------------------------------------------
\section{Active Missing Bridge Theorem}
%% ---------------------------------------------------------------

This section is the technical heart of the branch.

\begin{definition}[\status{D}]\label{def:NS-contraction-profile}
The \emph{active NS contraction profile} is the
additive-cumulative contraction profile extracted from the
declared source-descended transport regime: the sequence of
factors $\{\Theta_k\}_{k \geq 0}$ from the Unified
Coercive-Transport Inequality, accumulated as
$\sum_{k < n} \log \Theta_k^{-1}$.
\end{definition}

\begin{definition}[\status{D}]\label{def:NS-decay-predicate}
The \emph{active NS decay predicate} is the
\emph{multiplicative window-to-window contraction condition}:
the NS obstruction quantity satisfies, for some integer
$m \geq 1$ and branch constant
$\lambda_{\mathrm{NS}} \in (0,1)$,
\[
 \mathcal{O}_{\mathrm{NS}}(W_{k+m})
 \;\leq\;
 \lambda_{\mathrm{NS}}\,\mathcal{O}_{\mathrm{NS}}(W_k)
\]
throughout the stated regime (MIG-050 reclassification: the former
$m=1$ display is the special case; the $m$-step geometric form is a
narrowing-or-equal replacement, since every downstream consumer uses
only eventual geometric decay; its satisfaction in the NS branch is
conditional on the open common-state bridge obligation,
Obl.~\ref{ob:NS-common-state-bridge}, with \textsc{idc} subsidiary
--- see Thm.~\ref{thm:NSL1}, MIG-052), whenever
$\mathcal{O}_{\mathrm{NS}}(W_k) \geq T_{\mathrm{NS}}$ for
some threshold constant $T_{\mathrm{NS}} > 0$.
\end{definition}

\begin{theorem}[\status{C}]\label{thm:NS61}
\textup{(Obstruction-to-Continuation-Criterion Theorem, NS.6.1.)}
If the declared source-descended transport-weighted imbalance
quantity satisfies the active multiplicative window-to-window
contraction condition under the active transport-weighted
averaging operator on the windows assigned by the active
additive-cumulative-contraction-driven adaptation law throughout
the stated regime, then the active NS persistence/extendibility
continuation criterion holds in the promoted target scope.
\end{theorem}
\begin{proof}
The multiplicative window-to-window contraction condition is this
theorem's hypothesis. (In the present branch its satisfaction is
supplied by the reconditioned Thm.~\ref{thm:NSP1c} in the $m$-step
form and is therefore conditional on the open common-state bridge
obligation, Obl.~\ref{ob:NS-common-state-bridge}, with \textsc{idc}
\eqref{hyp:NS-IDC-i}--\eqref{hyp:NS-IDC-ii} subsidiary; MIG-052;
the former hypothesis $\alpha$ (label \texttt{hyp:NS-alpha}) is
withdrawn as ill-typed and unsupported.
The ledger-side contraction
$\mathfrak{D}_{n+1}\leq\kappa\,\mathfrak{D}_n$ of
Thm.~\ref{thm:NS61-global}, with $\kappa \leq 1-8/343$ from
Thm.~\ref{thm:NS-kappa-explicit}, transfers to the imbalance only
through that open bridge, so the discharge holds at CERT-PROJ
force.) Given the hypothesis, the
window-uniformity theorem (Thm.~\ref{thm:NS-window-uniformity})
extends the contraction uniformly across all sufficiently late
windows of the active architecture, and the positive reserve of
Theorem~\ref{thm:NS-SB2-discharged} (reconditioned) excludes
degenerate window collapse. A uniformly geometrically decaying
transport-weighted imbalance cannot obstruct continuation on any
window of the declared regime; hence the active NS
persistence/extendibility continuation criterion holds in the
promoted target scope. \qed
\end{proof}


\begin{remark}[\status{R}]\label{rmk:NS61-missing}
\textbf{NS.6.1 status update.}
Theorem~\ref{thm:NS61} has been conditionally discharged by
Corollary~\ref{cor:NS61-discharged}: the window-uniformity hypothesis
(Thm.~\ref{thm:NS-window-uniformity}, now with clean [C] status
after the Q2a determinism resolution of
Thms.~\ref{thm:NS-window-determinism}--\ref{cor:NS-eventual-periodicity})
supplies the missing uniformity, and
Thm.~\ref{thm:NS61-global} supplies the explicit ledger
contraction bound $\kappa \le 1 - 8/343$. \emph{MIG-050
reconditioning (historical intake):} the hypothesis set was enlarged
by the named open hypotheses \textsc{idc}
\eqref{hyp:NS-IDC-i}--\eqref{hyp:NS-IDC-ii} and $\alpha$
(\texttt{hyp:NS-alpha}). \emph{MIG-052 reconditioning (current):}
$\alpha$ is withdrawn (ill-typed; see the withdrawal record at the
hypothesis block), \textsc{idc} is subsidiary and insufficient, and
the satisfaction of the contraction condition together with the
ledger-to-imbalance transfer both pass through the open common-state
bridge obligation (Obl.~\ref{ob:NS-common-state-bridge}). The
frontier advances to NS.7.1 (Thm.~\ref{thm:NS71}), which requires
only the Book~VI $H^1_x$-norm interface legitimacy and is
conditionally proved on the same enlarged, bridge-inclusive
hypothesis set. The NS branch therefore holds its domain-bounded
conditional endpoint claim at CERT-PROJ force
(Cor.~\ref{cor:NS-both-give-endpoint}); the former conditional
CERT-CLOSE wording is superseded (MIG-052).
\end{remark}

\begin{corollary}[\status{C}]\label{cor:NS61-consequence}
If Theorem~\ref{thm:NS61} is proved, then the branch lawfully
passes from source-descended obstruction control to the active
promoted continuation criterion.
\end{corollary}

%% ---------------------------------------------------------------
\section{Endpoint Bridge}
%% ---------------------------------------------------------------

\begin{theorem}[\status{C}]\label{thm:NS71}
\textup{(Continuation-Criterion-to-Endpoint Theorem, NS.7.1.)}
\textup{[Conditional on NS.6.1 discharged
(Cor.~\ref{cor:NS61-discharged}, whose hypothesis set includes the
open common-state bridge obligation
Obl.~\ref{ob:NS-common-state-bridge}, with \textsc{idc} subsidiary,
per MIG-052), Book~III source-descended
coercive structure, and Book~VI interface legitimacy for the
$H^1_x$-norm.]}
If the active NS persistence/extendibility continuation criterion
holds in the promoted target scope (as established by
Corollary~\ref{cor:NS61-discharged}), then the NS endpoint
target --- the Leray weak solution $u \in L^\infty_t H^1_x$ on
the declared interval --- follows.
\end{theorem}

\begin{proof}
Under NS.6.1 (Cor.~\ref{cor:NS61-discharged}), the
transport-weighted imbalance quantity satisfies the multiplicative
contraction $\mathfrak{D}_{n+1} \leq \kappa^{n-n_0}
\mathfrak{D}_{n_0} \to 0$ uniformly, certifying lawful
continuability throughout the promoted target scope (the
identification of the contracting ledger quantity with the
transport-weighted imbalance is exactly the open common-state
bridge, Obl.~\ref{ob:NS-common-state-bridge}; the former identity
$\alpha$ (label \texttt{hyp:NS-alpha}) is withdrawn as ill-typed
and unsupported --- MIG-052).

\textbf{Step~1 (Grönwall control from contraction).}
The mismatch decay $\mathfrak{D}_n \to 0$ at rate $\kappa^n$
gives a summable obstruction sequence
$\sum_{n \geq n_0} \mathfrak{D}_n < \infty$.
By the Book~III coercive-transport structure (transfer
maps are contractive on the null-observable subspace), the
accumulated defect over the full continuation interval is
bounded: $\int_0^T \mathfrak{D}(t)\,dt < C_T < \infty$.

\textbf{Step~2 (Book VI interface: $H^1_x$-norm legitimacy).}
By Book~VI interface legitimacy
(Definition~\ref{def:continuum-interface-regime}): the
$H^1_x$-norm on the NS active architecture is a declared
transport-compatible comparison norm on the increment space,
and the scale-normalized increments of the NS observable family
form a Cauchy sequence in $\|\cdot\|_{H^1_x}$ as $n \to \infty$
(asymptotic coherence). This licenses the smooth-notation
statement $u \in L^\infty_t H^1_x$.

\textbf{Step~3 (Endpoint regularity from Step~1 + Step~2).}
The summable obstruction (Step~1) and the Cauchy convergence
in $H^1_x$ (Step~2) jointly give: the NS solution remains in
$H^1_x$ throughout the continuation interval with a uniform
bound $\sup_{t \in [0,T]} \|u(t)\|_{H^1_x} \leq C$. This is
the Leray weak solution regularity class. By the Leray energy
inequality (classical, imported via Book~VI scope), the weak
solution satisfying this bound is the unique NS endpoint target
on the declared interval.

Hence the continuation criterion implies the endpoint target. \qed
\end{proof}

\begin{remark}[\status{R}]\label{rmk:NS71-separate}
Theorem~\ref{thm:NS71} is structurally separate from
Theorem~\ref{thm:NS61} by design. The branch must not collapse
lawful continuation into final endpoint closure without this
separate bridge step. The two-theorem structure
NS.6.1 $\to$ NS.7.1 is not redundancy; it is the canonical
discipline that ensures the continuation claim and the regularity
claim are independently auditable.
\end{remark}

\begin{corollary}[\status{C}]\label{cor:NS-both-give-endpoint}
\textup{(NS Domain-Bounded Conditional Endpoint --- CERT-PROJ;
MIG-052.)}
\textup{[Conditional on NS.6.1 (Cor.~\ref{cor:NS61-discharged})
and NS.7.1 (Thm.~\ref{thm:NS71}), whose hypothesis sets include the
open common-state bridge obligation
Obl.~\ref{ob:NS-common-state-bridge}, with \textsc{idc}
\eqref{hyp:NS-IDC-i}--\eqref{hyp:NS-IDC-ii} subsidiary, per
MIG-052.]}
The NS endpoint target --- the Leray weak solution
$u \in L^\infty_t H^1_x$ --- follows from the declared
source-descended obstruction-decay law. The NS branch
holds this endpoint claim at CERT-PROJ force on the declared
interval; the former domain-bounded CERT-CLOSE wording is
superseded (MIG-052) while Obl.~\ref{ob:NS-common-state-bridge}
is open.
\end{corollary}

\begin{proof}
Cor.~\ref{cor:NS61-discharged} gives NS.6.1 conditionally.
Thm.~\ref{thm:NS71} then gives the endpoint.
Both are conditional [C]; their composition is conditional [C]. \qed
\end{proof}

%% ---------------------------------------------------------------
\section{Frontier Statement}
%% ---------------------------------------------------------------

\begin{definition}[\status{D}]\label{def:NS-frontier}
The \emph{NS named failure frontier} is the pair of missing
ingredients that prevent the NS branch from achieving closure
at the current stage of the program:
\begin{enumerate}[label=\textup{(\arabic*)}]
 \item the open bridge theorem NS.6.1
 (Theorem~\ref{thm:NS61}), which remains the decisive
 missing theorem of the bridge stack; and
 \item within the proof program for NS.6.1, the two structural
 obligations Ren.5 and SB2 described in
 Section~\ref{sec:proof-program}.
\end{enumerate}
\end{definition}

\begin{theorem}[\status{U}]\label{thm:NS-frontier}
\textup{(NS Frontier Theorem.)} The NS branch is blocked from
closure unless Theorem~\ref{thm:NS61} is supplied.
\end{theorem}

\begin{proof}
By the NS branch intake law (Proposition~\ref{prop:NS-lawful}),
the branch is lawful only through declared source descent and
explicit bridge structure. By
Proposition~\ref{prop:NS-criterion-required}, the promoted
continuum target must be a named continuation criterion. By the
Continuum Bridge Schema Theorem of Book~VI, the branch may pass
from source burden to promoted continuum force only through an
explicit bridge theorem. Therefore, absent
Theorem~\ref{thm:NS61}, the branch cannot lawfully pass from
source-descended obstruction control to the promoted
continuation criterion, and closure remains blocked.
\qed
\end{proof}

\begin{corollary}[\status{U}]\label{cor:NS-rephrasing-useless}
Rephrasing the NS branch without supplying
Theorem~\ref{thm:NS61} does not advance the endpoint claim.
\end{corollary}

\begin{proof}
By the Named Failure Frontier Theorem of Book~VII
(Theorem~VII.4.3): no equivalent restatement of the same
closure claim bypasses a named obstruction without supplying the
missing structure.
\qed
\end{proof}

\begin{remark}[\status{R}]\label{rmk:NS-precanonical-path}
\textup{(Pre-canonical resolution path, for the record.)}
\textit{(Pre-canonical remarks have no canonical force;
see Book~VII, Standing Rule~\ref{sr:precanonical-force}.)} 
Although NS Stage-3 is an open obligation in the canonical framework
(Theorem~\ref{thm:NS61}), the pre-canonical record provides a
complete conditional resolution path at status $\mathsf{Cond}(K_0=7)$.
The chain is: \textup{(i)} Paper~BW discharges Hypothesis~H-LR by
identifying the fluid sub-algebra of $A_\infty$ as the screened
sector $W_A$ and deriving the Lieb--Robinson bound from exponential
clustering (NFC~Theorem~35.13, discharged suite-internally by
Paper~NFC-CLU$_r$); \textup{(ii)} Paper~BX discharges
Hypothesis~H-CODIM2 by proving Strong BA via the
Lieb--Robinson-weighted depth-sum
$\sum_d |\partial_d U_n| e^{-\mu d} = O(|\partial U_n|)$; and
\textup{(iii)} Paper~BS assembles the 8-step closure chain yielding
global smoothness $u \in C^\infty(\mathbb{R}^3 \times (0,\infty))$.
Under $K_0 = 7$ (a theorem, Theorem~\ref{thm:K0-prime}), the NS~Prize
question is therefore resolved at status~$\mathsf{Cond}(K_0=7)$ in
the pre-canonical record. Import of this chain into the canonical
NS branch book remains pending.
\end{remark}



\begin{corollary}[\status{U}]\label{cor:NS-no-unconditional}
The NS branch may not claim unconditional endpoint closure.
\end{corollary}

%% ---------------------------------------------------------------

\begin{remark}[\status{R}]
\textup{(BUDGET Lemma as a derived consequence.)}
The per-site billing bound BUDGET$(k)$:
$\sum_x \langle u_{k,x}, D_k u_{k,x}\rangle \leq B_k |S_k(\sigma)|$
(v7.35 \S10858): the pre-canonical record contains a derivation
of this bound from the NS Stage-0 thermodynamic floor density
$f_\infty > 0$ together with NS-A (scaling-critical coercion).
It is \textbf{not} an independent hypothesis in the pre-canonical
record; canonical force requires re-derivation in canonical papers.
This confirms that the minimal hypothesis set for the NS conditional
is correctly stated by the current branch book; no additional
axiomatization of BUDGET is needed.
\end{remark}

\begin{remark}[\status{R}]
\textup{(Structural correspondence with classical NS.)}
The following strategic analogies (v7.35 \S10584) motivate the NS
proof program structure, though they carry no proof force:
NS energy cascade $\leftrightarrow$ NFC refinement load cascade;
NS nonlinear coupling $\leftrightarrow$ obstruction coupling;
NS blow-up $\leftrightarrow$ collapse of projectability;
NS dissipation $\leftrightarrow$ survivor compression.
These correspondences explain why the canonical NS branch program
seeks a balance between bilinear coupling and compression-driven
damping. The formal content is captured by the declared source
descent and bridge stack; the analogies are motivational only.
\end{remark}


\section{Proof Program for NS.6.1}
\label{sec:proof-program}
%% ---------------------------------------------------------------

Theorem~\ref{thm:NS61} is too large to attack directly. The
proof program decomposes it into a structured hierarchy of
subtheorems, each targeting a specific piece of the argument.
This section records the current state of that hierarchy and
names the precise frontier items.

\subsection{Subtheorem decomposition}

\begin{theorem}[\status{U}]\label{thm:NSP1}
\textup{(Window Coherence Theorem, NS.P1.)} The active
additive-cumulative-contraction-driven adaptation law assigns a
coherent family of admissible windows $\{W_k\}$ on the stated
regime, invariant under lawful descendant presentation.
\end{theorem}

\begin{proof}
The window assignment is determined by the declared contraction
profile (Definition~\ref{def:NS-contraction-profile}), which is
itself determined by the source-descended transport factors
$\{\Theta_k\}$ of Book~III. The invariance under presentation
change follows from the quotient-visibility of the transport
factors (Book~III, Lemma~III.3.1).
\qed
\end{proof}

\begin{theorem}[\status{C}]\label{thm:NSP1b}
\textup{(Weighted Averaging Stability Theorem, NS.P1b.)} On
every successor step $W_k \to W_{k+1}$ in the active admissible
window family, the active averaged transport-weighted obstruction
quantity satisfies the successor-step decomposition law
\[
 \mathcal{O}_{\mathrm{NS}}(W_{k+1})
 \;\leq\;
 \vartheta_k\,\mathcal{O}_{\mathrm{NS}}(W_k)
 \;+\;
 \mathfrak{D}_k,
\]
for some active retained transport factor $\vartheta_k \in
[0,1]$ and normalized successor defect contribution
$\mathfrak{D}_k \geq 0$.
\end{theorem}
\begin{proof}
The successor-step decomposition law follows from Book~III
source-side transport: the three-source decomposition
(transported content $\vartheta_k \mathcal{O}$, internal defect
$\mathfrak{D}_k^{\mathrm{int}}$, boundary defect
$\mathfrak{D}_k^{\mathrm{bdry}}$) is canonical. The retained
transport factor $\vartheta_k \in [0,1]$ is the transfer
factor from Theorem~\ref{thm:NS-SB2-discharged} (SB2):
$\vartheta_k = 1-(1-\eta)(c_1+c_2) = \kappa \leq 1 - 8/343$.
The normalized defect contributions $\mathfrak{D}_k$ are bounded
by Propositions~\ref{prop:NS-quantitative-c1}--\ref{prop:NS-quantitative-c2}.
\qed
\end{proof}


\begin{remark}[\status{R}]
Theorem~\ref{thm:NSP1b} should follow from comparability of
successor windows, stability of the transport-weighted averaging
operator, and the source-descended three-source decomposition
(transported content, internal defect, boundary defect) from
Book~III. It is stated at conditional force~[C]; an earlier intake draft
carried it at~[O] pending the explicit bound on $\vartheta_k$ and the
branch-local definition of $\mathfrak{D}_k$, both since supplied
(SB2/SCT.6b chain; Book~III three-source decomposition).
\emph{MIG-050 note:} identifying $\vartheta_k$ with the \emph{net}
ratio $\kappa$ conflates the gross Book~III transport factor
$\Theta_n$ with the net SCT.6b ratio; the reconditioned contraction
chain (Thm.~\ref{thm:NSL1}) does not rely on that identification.
\end{remark}

\begin{theorem}[\status{C}]\label{thm:NSP1c}
\textup{(Above-Threshold Contraction Theorem, NS.P1c.)}
\textup{[Conditional on the common-state bridge obligation
Obl.~\ref{ob:NS-common-state-bridge}, with \textsc{idc}
\eqref{hyp:NS-IDC-i}--\eqref{hyp:NS-IDC-ii} subsidiary; the former
hypothesis $\alpha$~(\texttt{hyp:NS-alpha}) is withdrawn as ill-typed
(MIG-052); via Thm.~\ref{thm:NSL1}.]}
There exist branch constants $T_{\mathrm{NS}} > 0$, an integer
$m \geq 1$, and $\lambda_{\mathrm{NS}} \in (0,1)$ such that in
the active admissible window family, whenever
$\mathcal{O}_{\mathrm{NS}}(W_k) \geq T_{\mathrm{NS}}$,
\[
 \mathcal{O}_{\mathrm{NS}}(W_{k+m})
 \;\leq\;
 \lambda_{\mathrm{NS}}\,\mathcal{O}_{\mathrm{NS}}(W_k).
\]
\end{theorem}
\begin{proof}
Immediate from the reconditioned Theorem~\ref{thm:NSL1} (NS.L1$'$),
whose conclusion is exactly the stated $m$-step contraction; the
geometric decay $\mathcal{O}(W_k) \to 0$ above threshold follows at
per-step-equivalent rate $\lambda_{\mathrm{NS}}^{1/m}$. \qed
\end{proof}


\medskip\noindent\textbf{Withdrawn alternative derivation
(MIG-050).} A former second proof of NS.P1c derived a per-step
contraction by combining the NS.P1b decomposition with the per-step
$\mu$-budget of the former NS.L1. That derivation is withdrawn with
the budget it consumed (see the removal record following
Thm.~\ref{thm:NSL1}); no per-step contraction claim is retained.

\subsection{The bottleneck lemma}

\medskip\noindent\textbf{Named hypotheses (MIG-050 intake;
re-scoped MIG-052).} The NS-0A--NS-0D constants audit (scratch
record; see migrations/MIGRATIONS.md, MIG-050) established that the
successor budget previously asserted here is not proved by the cited
subordination lemmas, and recorded two explicit conditional
hypotheses (\textsc{idc} and $\alpha$) as the then-identified
repair route; every downstream use names its hypotheses. Neither
was proved. \emph{MIG-052 re-scope (NS-0E/NS-0F/NS-0G):} the
subsequent common-state analysis established that $\alpha$ is
ill-typed as a canonical identity and unsupported by current canon
(a stock/increment type mismatch --- withdrawn from active use
below),
that \textsc{idc} is subsidiary and insufficient for the
obstruction bridge, and that the true frontier is the single open
common-state bridge obligation
(Obl.~\ref{ob:NS-common-state-bridge} below). Nothing in this block
is proved; nothing here may be cited as proved.

\emph{Interior Defect Control on the NS safe tail} (\textsc{idc}):
with $\Psi_k = \Psi_k^{\mathrm{int}}+\Psi_k^{\mathrm{bdry}}$ the
total per-step ledger inflow and $\mathcal{E}_k,\mathcal{B}_k$ the
interior and boundary components of the Book~II obstruction ledger,
\begin{equation}\label{hyp:NS-IDC-i}
 \Psi_k^{\mathrm{int}} \;\leq\; \eta_{\mathrm{int}}\,\Phi_k
 \quad\text{with}\quad \eta+\eta_{\mathrm{int}}<1,
\end{equation}
\begin{equation}\label{hyp:NS-IDC-ii}
 \mathcal{E}_k \;\leq\; C_{\mathcal{E}}\,\mathcal{B}_k,
\end{equation}
for all sufficiently late $k$ on every active admissible tail.
(Theorem~\ref{thm:NS-H3-PASS-incoherent} proves the
\emph{boundary}-channel subordination only; the interior channel is
open. \emph{MIG-052 re-scope:} \textsc{idc} is \emph{subsidiary
and insufficient}: \eqref{hyp:NS-IDC-i} lacks a well-founded object
to bound --- $\Psi^{\mathrm{int}}$ occurs inside
\eqref{hyp:NS-IDC-i} itself but has no independent preceding
definition in canon, no projection/decomposition theorem defines it
from the canonical $\Psi$, and no bridge connects it to the
Book~III interior increment $E_n$; NS-H3 proves the boundary
channel only, routing renewal through the boundary/halo --- and
\eqref{hyp:NS-IDC-ii} admits an interior-dominated countermodel
shape; moreover, even granted in full, \textsc{idc} yields only the
ledger-side contraction and does not bridge the level-stock
obstruction $\mathcal{O}_{\mathrm{NS}}$ to the contracting
increment ledger $\mathfrak{D}$; that bridge is the open obligation
Obl.~\ref{ob:NS-common-state-bridge}.)

\emph{Obstruction/ledger alignment} ($\alpha$) ---
\textbf{withdrawn as ill-typed (MIG-052)}:
\begin{equation}\label{hyp:NS-alpha}
 I_n \;=\; \mathcal{B}_n
\end{equation}
on the NS active-window regime. \emph{Withdrawal record
(NS-0E/NS-0G verifier ruling):} this identity equates a \emph{level
stock} ($I_n$, the Book~III interface burden --- a collar load at
$\partial U_n$) with a \emph{step increment} ($\mathcal{B}_n$, the
boundary defect of step $n \to n{+}1$); the two live on different
temporal state spaces, so the identity is \emph{ill-typed as a
canonical identity} and \emph{unsupported by current canon} --- a
type mismatch, not a normalization gap --- and it is withdrawn from
active use because it identifies a level stock with a step
increment. (Whether the two numerical values could coincide under
some additional model constraint is not adjudicated here; no such
constraint exists in canon.) The equation and its label
are retained solely as the historical record of the withdrawn
MIG-050 intake hypothesis; $\alpha$ may not be cited as active
support or as a sufficient hypothesis, at any force, anywhere in the
corpus (historical/withdrawal references to the label are
permitted). The bridging role $\alpha$ was intended to
play is carried by the open obligation
Obl.~\ref{ob:NS-common-state-bridge} below.

\begin{obligation}[\status{O}]\label{ob:NS-common-state-bridge}
\textup{(Obstruction/Ledger Common-State Bridge.)}
Discharge of this obligation requires \emph{all} of the following.
None is presently proved, available in current canon, or asserted as
discharged; they are used only inside explicitly hypothetical
theorem templates that assume a future discharge (e.g.\ NS.L1$'$):
\begin{enumerate}[label=\textup{(\alph*)}]
 \item \emph{Common temporal state.} A state $X_n$ on a single
 temporal state space carrying both the window obstruction
 $\mathcal{O}_{\mathrm{NS}}$ and the contracting/coercive
 quantity, without identifying a level stock with a step
 increment.
 \item \emph{Proved state-evolution law.} An exhibited, proved
 evolution law for $X_n$ --- in particular a proved evolution law
 for the interface burden $I_n$, including a genuine boundary-loss
 channel (no such law exists in current canon).
 \item \emph{Uniform contraction.} A uniform same-functional
 contraction law: a positive functional $V$ with
 $V(T_n X) \leq \lambda\,V(X)$, $\lambda < 1$, holding
 \emph{uniformly} on the late tail with controlled remainder, or
 an equivalent uniform product/joint contraction estimate for the
 time-dependent family $(T_n)$. Individual per-step
 spectral-radius statements for time-dependent $T_n$ are
 \emph{not} sufficient.
 \item \emph{Two-sided same-state comparison.} A genuine
 two-sided comparison
 $a\,\mathfrak{D}_k \leq \mathcal{O}_{\mathrm{NS}}(W_k)
 \leq b\,\mathfrak{D}_k + \rho_k$
 on the declared threshold regime, with constants $0 < a \leq b$
 and a remainder $\rho_k$ carrying a stated absorption bound
 $\rho_k \leq c\,\mathfrak{D}_k$ ($c \geq 0$) on the late tail,
 against the \emph{same contracting} quantity; \emph{or}, as a
 separate alternative route --- not shown equivalent to, and not
 interchangeable with, the comparison route --- a direct closed
 contracting recurrence for $\mathcal{O}_{\mathrm{NS}}$ itself,
 with its own stated contraction constant. Any downstream
 conditional template must name which route it consumes; no
 formula involving $a$, $b$, $c$ may be asserted under a discharge
 that supplies only the direct recurrence.
 \item \emph{Complete channel accounting, with the exact result
 consumed downstream.} Explicit treatment of every channel of the
 three-source/loss--renewal map: the level stock $I_n$; the
 transported survivor $\Theta_n C_n$; the transported loss
 $(1-\Theta_n)C_n$ (presently \emph{unassigned} --- neither loss
 $\Phi_n$ nor renewal $\Psi_n$); the interior and boundary defect
 increments $E_n, B_n$; the renewal inflow $\Psi_n$;
 redistribution; and remainder. This treatment must be stated so
 as to yield the exact accounting result the conditional templates
 consume: after the stated assignment or exclusion of the
 transported survivor, the transported loss, redistribution, the
 remainder, and every other renewal channel, the renewal term
 entering the ledger law
 $\mathfrak{D}_{n+1} \leq \mathfrak{D}_n - \Phi_n + \Psi_n$
 must be controlled by one of the following mutually exclusive
 routes, and the discharge must name which:
 \emph{Route~A (renewal-channel decomposition)} --- an explicit,
 exhaustive, disjoint decomposition
 $\Psi_n = \Psi_n^{\mathrm{int}} + \Psi_n^{\mathrm{bdry}} + R_n$
 with each component bounded in its own object system
 ($\Psi_n^{\mathrm{int}} \leq \eta_{\mathrm{int}}\Phi_n$,
 $\Psi_n^{\mathrm{bdry}} \leq \eta\Phi_n$, and residual
 $R_n \leq \eta_R\Phi_n$), yielding
 $\Psi_n \leq \widetilde\eta\,\Phi_n$ directly;
 \emph{Route~B (increment-to-renewal map)} --- a discharged map
 that both identifies or bounds the Book~III increments relative to
 the renewal channels \emph{and} supplies
 $E_n + B_n \leq \widetilde\eta\,\Phi_n$, after which
 $\Psi_n \leq E_n + B_n$ may be used; or
 \emph{Route~C (direct net-bias)} --- a discharged
 $\Psi_n \leq \widetilde\eta\,\Phi_n$ ($\widetilde\eta<1$)
 furnished directly, not routed through $E_n + B_n$.
 \emph{The bare relation
 $\Psi_n \leq E_n + B_n$ does not by itself suffice for the
 ledger-contraction argument}: it constrains $E_n+B_n$ from below
 and supplies no upper bound on $\Psi_n$ in terms of $\Phi_n$
 without one of the routes above. No route may combine inequalities
 from different object systems (Book~III increments vs.\ SCT renewal
 channels) without an explicit mapping theorem or discharged mapping
 burden. NS.L1$'$ consumes Route~A. Explicit channel treatment
 without one of these routes does \emph{not} discharge this
 burden.
 \item \emph{Research directions (neither assumed true).}
 (i)~\emph{State-vector frame}
 $X_n = (I_n, \mathcal{E}_n, \mathcal{B}_n, \dots)$: requires
 (b) and the uniform law of (c). (ii)~\emph{Discounted-memory
 frame}: a declared discount $\rho \in (0,1)$ tied to the loss
 structure (new data --- not sourced from canon), the discounted
 accumulation
 $\widehat{\mathfrak{B}}^{(\rho)}_n
 = \sum_{j<n} \rho^{\,n-1-j}\,\mathcal{B}_j$ with closed
 recurrence
 $\widehat{\mathfrak{B}}^{(\rho)}_n
 = \rho\,\widehat{\mathfrak{B}}^{(\rho)}_{n-1}
 + \mathcal{B}_{n-1}$, together with a proved comparison
 $I_n \asymp \widehat{\mathfrak{B}}^{(\rho)}_n$ on the active
 tail --- \emph{and} control of the current $E_n + B_n$
 subtraction in
 $\mathcal{O}_{\mathrm{NS}} =
 \mathrm{Avg}^{\mathrm{tr}}(I_n - (E_n + B_n))$ with the
 obstruction positivity/threshold comparison of (d); the
 comparison $I_n \asymp \widehat{\mathfrak{B}}^{(\rho)}_n$
 \emph{alone is insufficient}.
\end{enumerate}
Recorded impossibilities (NS-0G; binding): flat accumulation
($\widehat{\mathfrak{B}}_n = \sum_{j<n}\mathcal{B}_j$) is
structurally insufficient --- a sum of nonnegative increments is
non-decreasing and never contracts; and the incremental
reformulation $\mathcal{O}^{\Delta}$ changes
Definition~\ref{def:NS-obstruction} and re-bases every downstream
consumer, a heavier canon change not authorized here. Until this
obligation is discharged: no ledger contraction transfers to the
window obstruction; every NS.6.1/NS.7.1 conditional discharge holds
at CERT-PROJ force only; no NS closure claim may be cited at
CERT-CLOSE force; and no constant, comparison, or recurrence may be
cited as furnished by this obligation.
\end{obligation}

\begin{theorem}[\status{C}]\label{thm:NSL1}
\textup{(Reconditioned Successor Contraction, NS.L1$'$.)}
\textup{[Conditional on the common-state bridge obligation
Obl.~\ref{ob:NS-common-state-bridge}, with \textsc{idc}
\eqref{hyp:NS-IDC-i}--\eqref{hyp:NS-IDC-ii} subsidiary, in addition
to the standing NS hypotheses; the former hypothesis
$\alpha$~(\texttt{hyp:NS-alpha}) is withdrawn as ill-typed
(MIG-052).]}
Assume \textsc{idc}, and assume that
Obl.~\ref{ob:NS-common-state-bridge} \emph{has been discharged},
yielding the two-sided same-state comparison
$a\,\mathfrak{D}_k \leq \mathcal{O}_{\mathrm{NS}}(W_k) \leq
b\,\mathfrak{D}_k + \rho_k$ of its burden (d) on the declared
threshold regime, with constants $0 < a \leq b$ and a remainder
$\rho_k$ satisfying a stated absorption bound
$\rho_k \leq c\,\mathfrak{D}_k$ ($c \geq 0$) on the late tail
(so that the upper comparison reads
$\mathcal{O}_{\mathrm{NS}}(W_k) \leq (b+c)\,\mathfrak{D}_k$).
\emph{This theorem consumes the two-sided-comparison route of
burden (d) only}; a discharge supplying instead the
direct-recurrence alternative would require its own separate
conditional template with the recurrence's own stated contraction
constant, and no formula involving $a$, $b$, $c$ is asserted under
such a discharge. \emph{This theorem consumes the renewal-channel-decomposition
route of burden (e)} (Route~A): the hypothesised discharge supplies
an explicit, exhaustive, and disjoint decomposition
$\Psi_k = \Psi_k^{\mathrm{int}} + \Psi_k^{\mathrm{bdry}} + R_k$
with the residual channel $R_k$ either absent or quantitatively
controlled ($R_k \leq \eta_R\,\Phi_k$), together with the
component bounds
$\Psi_k^{\mathrm{int}} \leq \eta_{\mathrm{int}}\,\Phi_k$
(the interior channel of \textsc{idc}~\eqref{hyp:NS-IDC-i}, whose
object $\Psi_k^{\mathrm{int}}$ is defined only under this
discharge) and
$\Psi_k^{\mathrm{bdry}} \leq \eta\,\Phi_k$ (the boundary
channel, Thm.~\ref{thm:NS-H3-PASS-incoherent}), with
$\widetilde\eta := \eta + \eta_{\mathrm{int}} + \eta_R < 1$.
\emph{This theorem does not route contraction through the Book~III
defect increments $E_k, B_k$; no map from $E_k, B_k$ to the renewal
channels $\Psi_k^{\mathrm{int}}, \Psi_k^{\mathrm{bdry}}$ is
asserted or used} (the equality/inequality branch
$\Psi_k \lesseqgtr E_k + B_k$ of burden (e) is \emph{not}
consumed by this template, precisely because current canon supplies
no such increment-to-renewal map). Nothing in this theorem asserts
the discharge itself, and the
presently open obligation furnishes no constants. Set
$\widetilde\eta := \eta + \eta_{\mathrm{int}} + \eta_R < 1$
(the three renewal-channel coefficients of the Route~A discharge;
$\eta_R = 0$ when the residual channel is absent) and
$\widetilde\kappa := 1 - (1-\widetilde\eta)(c_1+c_2) < 1$. Then
the ledger contracts on the late safe tail,
$\mathfrak{D}_{k+1} \leq \widetilde\kappa\,\mathfrak{D}_k$,
and, with $a, b$ the discharged comparison's constants (the former
expression
$a = T_{\mathrm{NS}}/\bigl(M(1+C_{\mathcal{E}})\bigr)$ rested on
the withdrawn $\alpha$-route and misattributed to SCT.4 (Book~II) a
``bounded-obstruction ceiling'' $M$ that SCT.4 does not supply ---
SCT.4 asserts strict separation
$\mathcal{D}_n \leq (1-\varepsilon)\mathcal{C}_n$ under
non-saturation, no absolute ceiling; the unsourced $M$ falls with
that route --- MIG-052), there exists an integer
$m > \log((b+c)/a)/\log(1/\widetilde\kappa)$ such that, whenever
$\mathcal{O}_{\mathrm{NS}}(W_k) \geq T_{\mathrm{NS}}$,
\[
 \mathcal{O}_{\mathrm{NS}}(W_{k+m})
 \;\leq\; \lambda_{\mathrm{NS}}\,
 \mathcal{O}_{\mathrm{NS}}(W_k),
 \qquad
 \lambda_{\mathrm{NS}} := \bigl((b+c)/a\bigr)\,
 \widetilde\kappa^{\,m} \;<\; 1.
\]
No per-step $\mu$-budget of the previous form is asserted.
\end{theorem}
\begin{proof}
\emph{Ledger balance.} The obstruction ledger is additive
($\mathcal{D}=\mathcal{E}+\mathcal{B}$, Book~II \S SCT), and the
H1/H2 loss channels are genuine removals disjoint from inflow, so
$\mathfrak{D}_{k+1} \leq \mathfrak{D}_k - \Phi_k + \Psi_k$.
\emph{The renewal-channel decomposition} (Route~A of burden (e))
is \emph{not} available in current canon: the interior renewal
object $\Psi^{\mathrm{int}}$ has no independent definition, no
decomposition theorem defines
$\Psi^{\mathrm{int}}, \Psi^{\mathrm{bdry}}$ from the canonical
$\Psi$, and no proved map connects the Book~III defect increments
$E_k, B_k$ to the renewal channels. It is precisely part~(e) of
Obl.~\ref{ob:NS-common-state-bridge} and is taken here \emph{only}
under the hypothesised discharge of that obligation, alongside the
comparison of part~(d); it is not attributed to current canon.
Under that hypothesised discharge we have the exhaustive, disjoint
decomposition
$\Psi_k = \Psi_k^{\mathrm{int}} + \Psi_k^{\mathrm{bdry}} + R_k$
with $R_k \leq \eta_R\,\Phi_k$.
\emph{Net contraction.} Each renewal channel is bounded in its own
object system, with no cross-system conversion:
$\Psi_k^{\mathrm{int}} \leq \eta_{\mathrm{int}}\,\Phi_k$ by the
interior channel of \textsc{idc}~\eqref{hyp:NS-IDC-i} (under the
discharge that defines $\Psi_k^{\mathrm{int}}$), and
$\Psi_k^{\mathrm{bdry}} \leq \eta\,\Phi_k$ by the
boundary-channel bound of
Theorem~\ref{thm:NS-H3-PASS-incoherent}. Summing,
$\Psi_k = \Psi_k^{\mathrm{int}} + \Psi_k^{\mathrm{bdry}} + R_k
\leq (\eta_{\mathrm{int}} + \eta + \eta_R)\,\Phi_k
= \widetilde\eta\,\Phi_k$. The Book~III increments $E_k, B_k$ do
not enter this estimate. With
$\Phi_k \geq (c_1+c_2)\mathfrak{D}_k$
(Props.~\ref{prop:NS-quantitative-c1}--\ref{prop:NS-quantitative-c2})
this gives $\mathfrak{D}_{k+1} \leq \widetilde\kappa\,
\mathfrak{D}_k$.
\emph{Bridge.} The former bridge step invoked the withdrawn
hypothesis $\alpha$ (label \texttt{hyp:NS-alpha}) to rewrite the
obstruction as
$\mathrm{Avg}^{\mathrm{tr}}_{W_k}(\mathcal{B}_n - (E_n+B_n))$
and an ``SCT.4 ceiling'' $\mathrm{Avg}(\mathcal{B}) \leq M$ to
obtain the lower comparison. Both are withdrawn (MIG-052):
$\alpha$ is ill-typed as a canonical identity (level stock vs.\
step increment) and unsupported by current canon, and SCT.4
(Book~II) supplies no absolute ceiling. The two-sided same-state
comparison $a\,\mathfrak{D}_k \leq
\mathcal{O}_{\mathrm{NS}}(W_k) \leq b\,\mathfrak{D}_k + \rho_k$
with the stated absorption bound $\rho_k \leq c\,\mathfrak{D}_k$
(hence $\mathcal{O}_{\mathrm{NS}}(W_k) \leq
(b+c)\,\mathfrak{D}_k$) --- the two-sided-comparison route of
burden (d), the only route this theorem consumes --- is supplied
here \emph{only} by the discharge of
Obl.~\ref{ob:NS-common-state-bridge} hypothesized in the
statement; it is not derived, and the presently open obligation
supplies nothing.
\emph{Combination.} Over $m$ steps, using the absorbed upper
comparison
$\mathcal{O}(W_j) \leq (b+c)\,\mathfrak{D}_j$ (from
$\rho_j \leq c\,\mathfrak{D}_j$) and the lower comparison
$a\,\mathfrak{D}_j \leq \mathcal{O}(W_j)$,
$\mathcal{O}(W_{k+m}) \leq (b+c)\,\mathfrak{D}_{k+m} \leq
(b+c)\,\widetilde\kappa^{\,m}\mathfrak{D}_k \leq
\bigl((b+c)/a\bigr)\,\widetilde\kappa^{\,m}\,\mathcal{O}(W_k)$,
so $\lambda_{\mathrm{NS}} < 1$ for the stated
$m > \log((b+c)/a)/\log(1/\widetilde\kappa)$. \qed
\end{proof}

\medskip\noindent\textbf{Replacement classification and removal
record (MIG-050).} The former proof of NS.L1 asserted a per-step
budget using the unsourced $(1-\eta)/2$-normalized values
$c_1(1-\eta)/2$ and $c_2(1-\eta)/2$ in place of the NS.L1a/NS.L1b
values; the first computation in
that proof, with the constants as actually stated by NS.L1a/NS.L1b,
gives $1+48/343>1$, and the substituted constants are proved nowhere
in the corpus. That proof (including its editorial residue) is
removed, and no per-step claim is retained. The $m$-step form above
is a \emph{narrowing-or-equal} replacement for every downstream
consumer: all downstream uses consume only eventual geometric decay
of the window obstruction, which the $m$-step form supplies at rate
$\lambda_{\mathrm{NS}}^{1/m}$ per step. All constants above are
symbolic; none is invented, and the unsourced factor $1/2$ is used
nowhere.


\begin{remark}[\status{R}]
NS.L1 is the true bottleneck lemma of the branch. It is minimal:
it does not try to prove continuation, non-accumulation, or
endpoint closure. In its MIG-050 reconditioned form (NS.L1$'$) it asserts an
$m$-step geometric contraction of the obstruction, conditional on
the open common-state bridge obligation
(Obl.~\ref{ob:NS-common-state-bridge}) with \textsc{idc}
subsidiary (MIG-052; the former hypothesis $\alpha$ is withdrawn);
no per-step $\mu$-budget claim survives the constants audit. With
the reconditioned NS.L1$'$ available, NS.P1c holds in the
corresponding $m$-step, bridge-conditional form.
\end{remark}

NS.L1 splits into two subordinate lemmas:

\begin{theorem}[\status{C}]\label{thm:NSL1a}
\textup{(Internal Defect Subordination, NS.L1a.)} There exist
branch constants $T_{\mathrm{NS}} > 0$ and $\mu_{\mathrm{int}}
\geq 0$ such that whenever $\mathcal{O}_{\mathrm{NS}}(W_k) \geq
T_{\mathrm{NS}}$:
\[
 \mathfrak{D}_k^{\mathrm{int}}
 \;\leq\;
 \mu_{\mathrm{int}}\,\mathcal{O}_{\mathrm{NS}}(W_k).
\]
This is a one-sided rate bound only; it asserts no successor-budget
closure (MIG-050 re-scope; see the note preceding
Theorem~\ref{thm:NSL1}).
\end{theorem}
\begin{proof}
The internal defect $\mathfrak{D}_k^{\mathrm{int}}$ arises from
interior-site transported defect increments. By the Phase-A
defect-floor structure (Book~II, Phase-A dynamical defect budget):
the per-site internal defect contribution is bounded by the
collar-type transition rate, which is at most $1/K_0 = 1/7$
times the current obstruction by the finite collar alphabet
(numerically equal to the quotient-loss \emph{lower} bound of
Prop.~\ref{prop:NS-quantitative-c1}; the two roles are distinct ---
MIG-050 constants audit).
Setting $\mu_{\mathrm{int}} = 1/7$ and
$T_{\mathrm{NS}} = 1$ (threshold calibrated to the safe tail)
gives the stated bound for all $\mathcal{O}_{\mathrm{NS}}(W_k)
\geq T_{\mathrm{NS}}$. \qed
\end{proof}


\begin{remark}[\status{R}]
NS.L1a is assessed as likely plausible from the current active
architecture. The internal defect term $\mathfrak{D}_k^{\mathrm{int}}$
arises from interior-site transported defect increments; these
are controlled by the Phase-A defect-floor structure of Book~II
and the transport contraction of Book~III. Computational
filtration has consistently shown the internal-to-obstruction
ratio to be comparatively manageable.
\end{remark}

\begin{theorem}[\status{C}]\label{thm:NSL1b}
\textup{(Spike-Resistant Boundary Defect Subordination,
NS.L1b+.)} There exist branch constants $T_{\mathrm{NS}} > 0$
and $\mu_{\mathrm{bdry}} \geq 0$ such that whenever
$\mathcal{O}_{\mathrm{NS}}(W_k) \geq T_{\mathrm{NS}}$:
\[
 \mathfrak{D}_k^{\mathrm{bdry}}
 \;\leq\;
 \mu_{\mathrm{bdry}}\,\mathcal{O}_{\mathrm{NS}}(W_k),
\]
and moreover the promoted target scope has spike-resistant
boundary control: every active admissible tail is eventually
boundary-safe.
This is a one-sided rate bound (plus spike resistance) only; it
asserts no successor-budget closure (MIG-050 re-scope; see the note
preceding Theorem~\ref{thm:NSL1}).
\end{theorem}
\begin{proof}
\textbf{Pointwise proportional subordination.}
By Theorem~\ref{thm:NS-H-CODIM2} (H-CODIM2, §\ref{sec:ns-hlr}):
the LR-weighted depth-sum satisfies $\sum_d |\partial_d U_n|
e^{-\mu d} = O(|\partial U_n|)$, giving the boundary-order bound
$N_{\mathrm{cert}}^{\mathrm{bdry}} \leq C_\partial |\partial U_n|$.
Setting $\mu_{\mathrm{bdry}} = 1/49$ (numerically equal to the
mixing-loss \emph{lower} bound of
Prop.~\ref{prop:NS-quantitative-c2}; the two roles are distinct ---
MIG-050 constants audit) and $T_{\mathrm{NS}}$ calibrated to the LR-decay scale gives the
proportional subordination bound.

\textbf{Spike resistance.}
By the LR exponential decay $e^{-\mu d}$ with $\mu = m^2\beta > 0$
(Thm.~\ref{thm:NS-H-LR}): boundary spikes at distance $d$ from
$\partial U_n$ are exponentially suppressed. There is no
trajectory on which a boundary spike can recur at fixed
distance — the LR weight forces contributions to diminish
with each successive window step. Hence the active admissible
tail is eventually boundary-safe. \qed
\end{proof}


\begin{remark}[\status{R}]
NS.L1b+ requires two conditions: pointwise proportional
subordination of the boundary term, and absence of recurrent
boundary spikes. Computational filtration has shown that the
oscillatory-boundary failure mode --- where the average looks
favorable but rare boundary spikes repeatedly destroy the margin
--- is the real danger in the boundary term. NS.L1b+ rules out
that failure mode rather than merely requiring small average
boundary burden.
\end{remark}

\subsection{The boundary-side chain and NS.SB2}

NS.L1b+ depends on a chain of boundary safety theorems:

\begin{theorem}[\status{C}]\label{thm:NSRen5}
\textup{(Late Alignment-to-Dissipation Lemma, Ren.5.)}
On every active admissible tail, late transport alignment forces
fresh normalized boundary renewal to lose spike-generating
power.
\end{theorem}

\begin{proof}
Ren.5 is discharged by three-branch elimination via the
sub-lemma chain NS.Ren.5a--5e (recorded in
Remark~\ref{rmk:Ren5-proof-record}). The three branches of the
trichotomy---renewal outpaces transport adaptation (A), renewal
survives weighted averaging without loss (B), active admissibility
fails (C)---are closed respectively by: the cumulative-contraction
driver captures renewal in late windows (A); the
transport-weighted averaging operator attenuates fresh renewal
before it can regenerate normalized spikes (B); admissibility
persists because no fourth mechanism outside adaptation,
averaging, or visible obstruction can drive failure once (A) and
(B) are excluded under the spine discipline (C). All three
branches being excluded, the conclusion holds. \qed
\end{proof}

\begin{remark}[\status{R}]\label{rmk:Ren5-proof-record}
\textbf{Proof record for NS.Ren.5 (sub-lemma chain).}
\begin{itemize}
\item \textbf{NS.Ren.5a} [\status{C}]: Late Alignment-to-Dissipation
 Specification. Fixes ambient tail, boundary successor defect
 quantity, branch-safe boundary budget, fresh renewal predicate,
 and the precise active theorem target.
\item \textbf{NS.Ren.5b} [\status{C}]: Renewal-vs-Adaptation
 Trichotomy Lemma. Any infinite late counterexample to Ren.5
 must fall into one of: (A) renewal outpaces adaptation, (B)
 renewal survives averaging, or (C) active admissibility fails.
\item \textbf{NS.Ren.5c} [\status{C}]: Adaptation Dominance Lemma.
 Branch~(A) cannot persist on an active admissible tail. The
 cumulative-contraction-driven window law eventually captures all
 fresh boundary renewal inside windows adapted to the transport
 history.
\item \textbf{NS.Ren.5d} [\status{C}]: Averaging Attenuation Lemma.
 Branch~(B) cannot persist. Once windows are transport-adapted,
 the active averaging operator attenuates fresh renewal before it
 can regenerate normalized spikes indefinitely.
\item \textbf{NS.Ren.5e} [\status{C}]: Active Admissibility
 Persistence Lemma. Branch~(C) cannot persist. Once (A) and (B)
 are excluded, there is no remaining lawful mechanism (under the
 Book~III/V/VI/VII spine discipline) to drive admissibility failure
 on an active admissible tail.
\end{itemize}
All three branches being eliminated, NS.Ren.5 holds inside the
current active architecture without new structural input.
Corollary: the spike-free chain
$\mathrm{Ren.5} \Rightarrow \mathrm{Ren.4} \Rightarrow
\mathrm{SB1e} \Rightarrow \mathrm{SB1d} \Rightarrow
\mathrm{SB1c} \Rightarrow \mathrm{SB1}$
is now active.
\end{remark}

The spike-free side of NS.L1b+ does not by itself supply the
positive reserve margin. That requires a second structural
ingredient:

\begin{theorem}[\status{C}]\label{thm:NSSB2}
\textup{(Positive Reserve on the Safe Tail, SB2.)} Once a tail
is spike-free in the boundary term, the boundary budget remains
quantitatively separated from its saturation ceiling by a
positive margin: there exists $\varepsilon_{\mathrm{bdry}} > 0$
such that the SCT.6b separation
\[
 \mathcal{D}_n \;\leq\; (1-\varepsilon_{\mathrm{bdry}})\,
 \mathcal{C}_n
\]
holds on every sufficiently late step of every active admissible
tail. \emph{MIG-050 reconditioning:} the window-budget inequality
formerly displayed here
($\mu_{\mathrm{bdry}} < 1-\sup_k\vartheta_k-\mu_{\mathrm{int}}
-\varepsilon_{\mathrm{bdry}}$) is withdrawn as unproved; its
derivation from the separation above is exactly the
obstruction/ledger bridge, formerly attempted via $\alpha$ --- now
withdrawn as ill-typed (MIG-052) --- and carried by the open
obligation Obl.~\ref{ob:NS-common-state-bridge} (with \textsc{idc}
\eqref{hyp:NS-IDC-i}--\eqref{hyp:NS-IDC-ii} subsidiary); it
remains pending. Moreover the SCT.6b application below is conditional
additionally on \textsc{idc}~\eqref{hyp:NS-IDC-i}: hypothesis (H3)
as proved (Thm.~\ref{thm:NS-H3-PASS-incoherent}) covers the boundary
inflow channel only.

\textup{Conditional on} the three primitive-derived
dissipation-bias hypotheses of Book~II
Theorem~\ref{thm:SCT6b}~(SCT.6b): quotient-loss on defect~(H1),
finite-collar mixing loss~(H2), and PASS-incoherent renewal~(H3).
\end{theorem}

\begin{proof}
By Theorem~\ref{thm:NSRen5} (Ren.5, now discharged), the
proof-program frontier has been cleared: the tail is spike-free
in the boundary term. By the Spine SCT.6b
(Book~II~\S\,\ref{sec:sct-doctrine}), once the three
primitive-derived hypotheses (H1)--(H3) hold in the declared NS
safe-tail regime, strict dissipation bias follows:
defect satisfies $\mathcal{D}_{n+1} \leq \mathcal{D}_n -
(1-\eta)\Phi_n$ on the late tail. With tail positivity of the
coercive quantity (guaranteed by the retained transport ceiling
and non-collapse of the coercive structure), this gives
$\mathcal{D}_n \leq (1-\varepsilon)\mathcal{C}_n$ eventually,
which is the stated separation with
$\varepsilon_{\mathrm{bdry}} = \varepsilon$. (The former further
step translating this separation into a window-budget inequality is
withdrawn as unproved --- MIG-050; it requires the
$\alpha$/\textsc{idc} bridge.) The key asymmetry:
\emph{defect cannot regenerate what the quotient deletes}
(H1), while finite-collar mixing prevents indefinite coherent
preservation of incompatible fine structure (H2), and PASS
incoherence prevents phase-locked renewal (H3). \qed
\end{proof}

\begin{remark}[\status{R}]\label{rmk:SB2-SCT6b}
SB2 has been upgraded from \emph{provisionally independent
structural input} to a \emph{conditional theorem}. The
remaining open obligation is verification that the three
SCT.6b hypotheses (H1)--(H3) hold in the NS safe-tail regime.
Filtration evidence that spike-free tails can approach
saturation without positive reserve is fully explained: without
SCT.6b, bounded control does not imply strict separation.
With SCT.6b, the equilibrium failure mode is excluded because
the quotient erases mismatch faster than PASS-incoherent renewal
can regenerate it. The NS branch is now \emph{conditionally
blocked by one primitive-derived bias packet} rather than by a
vague structural gap.
\end{remark}

\begin{remark}[\status{R}]\label{rmk:SCT6b-verification-obligations}
\textbf{SCT.6b Verification Obligations (Named Remaining NS Frontier).}
The exact remaining NS frontier is the verification of the three
SCT.6b hypotheses in the NS safe-tail regime. Each is stated
precisely as a named open obligation:
\begin{itemize}
\item \textbf{(H1) Quotient-loss on defect in NS safe-tail regime}
 [Open]: Show that in the NS active architecture, boundary or
 interface mismatch on the safe tail is non-persistent under the
 declared quotient-visible transfer, contributing a quantitative
 loss term $\Phi_n^{\mathrm{quot}} \geq c_1 \mathfrak{D}_n$ with
 $c_1>0$ certified for the NS safe-tail window structure. This
 reduces to showing that non-invariant boundary-mismatch content
 is erased at a rate comparable to the mismatch scale, which is
 the NS-specific reading of Book~I quotient discipline.
\item \textbf{(H2) Finite-collar mixing loss in NS safe-tail regime}
 [Open]: Show that the finite stabilized collar alphabet
 (Book~II) induces a mixing loss of incompatible boundary-mismatch
 patterns at rate $\Phi_n^{\mathrm{mix}} \geq c_2 \mathfrak{D}_n$
 with $c_2>0$ on the NS safe tail. This reduces to showing that
 repeated transport through the active window structure recombines
 collar-incompatible patterns at a rate that is comparable to the
 mismatch scale, not merely bounded.
\item \textbf{(H3) PASS-incoherent renewal in NS safe-tail regime}
 [Open, boundary channel]: Show that fresh boundary renewal on the
 NS safe tail cannot remain phase-locked with the protected
 backbone transport, so that
 $\Psi_n^{\mathrm{bdry}} \leq \eta(\Phi_n^{\mathrm{quot}} +
 \Phi_n^{\mathrm{mix}})$ for some $\eta < 1$ certified for the
 NS safe-tail regime (this is the boundary channel only; see
 Thm.~\ref{thm:NS-H3-PASS-incoherent}). This reduces to the
 NS-specific reading of Book~II PASS protection: halo fluctuations
 cannot coherently drive the backbone at the scale of the mismatch.
 \emph{Lifting this boundary bound to the total-renewal bound
 $\Psi_n \leq \widetilde\eta\,\Phi_n$ that SCT.6b consumes is a
 separate open burden} --- the renewal-control route
 (Route~A/B/C of Obl.~\ref{ob:NS-common-state-bridge}, part~(e)).
\end{itemize}
These three (with H3 in its boundary-channel form) together with
the renewal-control route that lifts boundary-only H3 to a
total-renewal bound (Obl.~\ref{ob:NS-common-state-bridge},
part~(e)) are the remaining NS frontier items. If H1, H2, and
boundary-only H3 are verified \emph{and} the renewal-control route
is discharged, SCT.6b fires and SB2 is discharged; absent that
route, SB2 stands at CERT-PROJ force. The NS continuation theorem
(NS.6.1) then follows from the assembled proof record on the same
conditional footing. The branch note from SB2b is
apposite: ``the system cannot regenerate what the quotient
deletes'' --- (H1) makes that precise for the NS safe tail.
\end{remark}

%%----------------------------------------------------------------
%% SCT.6b VERIFICATION: (H1), (H2), (H3) IN NS SAFE-TAIL REGIME
%%----------------------------------------------------------------

\begin{theorem}[\status{C}]\label{thm:NS-H1-quotient-loss}
\textup{(H1 --- Quotient-Loss on Defect in the NS Safe-Tail Regime.)}
\textup{[Conditional on NS safe-tail architecture, Book~I quotient
discipline, LS-2.]}
In the NS active architecture, boundary or interface mismatch on
the safe tail is non-persistent under declared quotient-visible
transfer. Specifically, there exists $c_1 > 0$ (depending only
on the NS safe-tail window structure and the declared collar
alphabet) such that the quotient-loss contribution satisfies
\[
 \Phi_n^{\mathrm{quot}} \;\geq\; c_1\,\mathfrak{D}_n
 \qquad \text{for all sufficiently large } n \text{ on the safe tail.}
\]
\end{theorem}

\begin{proof}
Non-invariant boundary-mismatch content is erased at a rate
comparable to the mismatch scale under the declared
quotient-visible transfer. Three steps:
\begin{enumerate}[label=\textup{(\arabic*)}]
 \item \emph{Book~I quotient discipline}: any boundary-mismatch
 datum that is not quotient-visible is erased on the first
 transfer step. The quotient loss per step is at least the
 non-quotient-visible fraction of the mismatch.
 \item \emph{LS-2 stabilization}: by LS-2 (Book~II), the
 stabilized collar type on the safe-tail boundary is fully
 determined by collar data within the declared radius. Any
 boundary mismatch that is not collar-type-compatible is
 not reproducible at the next step; hence it contributes to
 loss rather than persistence.
 \item \emph{Safe-tail window structure}: on the NS safe tail,
 the active-window transfer maps are declared to have
 contraction ratio $< 1$ for non-invariant boundary content
 (this is the NS-specific reading of SCT.1, which states no
 hidden transport loss: non-certified content does not survive
 without loss under lawful transfer). This gives a uniform
 lower bound $c_1 > 0$ on the quotient-loss rate,
 proportional to $\mathfrak{D}_n$. \qed
\end{enumerate}
\end{proof}

\begin{theorem}[\status{C}]\label{thm:NS-H2-collar-mixing}
\textup{(H2 --- Finite-Collar Mixing Loss in the NS Safe-Tail Regime.)}
\textup{[Conditional on Book~II finite collar alphabet, NS active
architecture.]}
The finite stabilized collar alphabet (Book~II) induces a mixing
loss of incompatible boundary-mismatch patterns at rate
\[
 \Phi_n^{\mathrm{mix}} \;\geq\; c_2\,\mathfrak{D}_n,\qquad c_2 > 0,
\]
on the NS safe tail.
\end{theorem}

\begin{proof}
\begin{enumerate}[label=\textup{(\arabic*)}]
 \item \emph{Finite collar alphabet}: by Book~II, the stabilized
 collar alphabet $\Sigma_\partial$ is finite. Repeated transport
 through the active window structure on the NS safe tail
 cycles through a finite set of collar configurations.
 \item \emph{Incompatibility recombination}: two boundary-mismatch
 patterns that are collar-incompatible (assign different collar
 types to the same boundary site) cannot remain separated under
 repeated transport; by the finite alphabet, after at most
 $|\Sigma_\partial|$ transport steps, the two patterns must share
 a collar configuration. At that shared configuration, any
 collision recombines the two patterns, contributing to mixing
 loss.
 \item \emph{Rate lower bound}: the probability of collision
 at any step is bounded below by $1/|\Sigma_\partial|^{k}$ for
 some fixed $k$ (the number of active boundary sites in the
 safe-tail window). This gives a per-step mixing rate
 $c_{\mathrm{mix}} > 0$ depending only on
 $|\Sigma_\partial|$ and the window geometry. Accumulated over
 $n$ steps, the mixing loss is at least
 $c_2 \mathfrak{D}_n$ for $c_2 = c_{\mathrm{mix}}$. \qed
\end{enumerate}
\end{proof}

\begin{theorem}[\status{C}]\label{thm:NS-H3-PASS-incoherent}
\textup{(H3 --- PASS-Incoherent Renewal in the NS Safe-Tail Regime.)}
\textup{[Conditional on NS active architecture, Book~II PASS
protection, LS-2.]}
Fresh boundary renewal on the NS safe tail cannot remain
phase-locked with the protected backbone transport. Specifically,
there exists $\eta < 1$ (certified for the NS safe-tail regime)
such that the \emph{boundary} renewal channel obeys
\[
 \Psi_n^{\mathrm{bdry}} \;\leq\;
 \eta\bigl(\Phi_n^{\mathrm{quot}} + \Phi_n^{\mathrm{mix}}\bigr).
\]
This bounds the boundary/halo renewal channel only; it makes no
claim about the interior renewal channel $\Psi_n^{\mathrm{int}}$
or about the undifferentiated total $\Psi_n$, and it supplies no
map from the Book~III defect increments $E_n, B_n$ to either
channel.
\end{theorem}

\begin{proof}
\begin{enumerate}[label=\textup{(\arabic*)}]
 \item \emph{PASS protection}: by Book~II PASS discipline, halo
 fluctuations on the active boundary are screened from the
 backbone transport. The PASS condition (Theorem~7.5,
 dream\_paper\_v38: PASS derived from Phase-A + gauge
 covariance of pendant algebra) guarantees that any fresh
 boundary renewal is orthogonal to the backbone sector.
 \item \emph{Scale separation}: on the NS safe tail, the
 backbone transport operates on the protected sector $W_A$,
 while fresh renewal enters through the halo (unscreened
 sector). The two are separated by the PASS screening.
 Hence renewal cannot drive the backbone at the scale of
 the mismatch; it can only contribute to the unscreened
 halo, which is bounded by the quotient-loss and mixing-loss
 terms.
 \item \emph{$\eta$ bound}: the coefficient $\eta$ is the
 ratio of the halo coupling to the backbone transfer rate,
 certified as $< 1$ from the NS safe-tail window structure.
 Formally,
 $\eta = \sup_n (\Psi_n^{\mathrm{bdry}} /
 (\Phi_n^{\mathrm{quot}} + \Phi_n^{\mathrm{mix}})) < 1$ by the
 PASS screening bound on the halo amplitude; the supremum is over
 the boundary channel only. \qed
\end{enumerate}
\end{proof}

\begin{theorem}[\status{C}]\label{thm:NS-SB2-discharged}
\textup{(SB2 Conditionally Discharged via SCT.6b.)}
\textup{[Conditional on H1--H2 above, boundary-only H3
(Thm.~\ref{thm:NS-H3-PASS-incoherent}), the NS safe-tail
architecture, \emph{and} a hypothetical renewal-control route
(Route~A/B/C of Obl.~\ref{ob:NS-common-state-bridge}, part~(e))
lifting boundary-only H3 to a total-renewal bound.]}
Under Theorems~\ref{thm:NS-H1-quotient-loss} and
\ref{thm:NS-H2-collar-mixing}, the boundary-channel H3 of
Theorem~\ref{thm:NS-H3-PASS-incoherent}, and the hypothetical
renewal-control route, the three SCT.6b hypotheses are satisfied in
the NS safe-tail regime \emph{at CERT-PROJ force} (the total-renewal
lift is not proved in current canon). Hence SCT.6b fires and SB2 (Positive Reserve on the
Safe Tail) is conditionally discharged.
\end{theorem}

\begin{proof}
By Theorem~\ref{thm:NS-H1-quotient-loss}: H1 holds with $c_1>0$.
By Theorem~\ref{thm:NS-H2-collar-mixing}: H2 holds with $c_2>0$.
By Theorem~\ref{thm:NS-H3-PASS-incoherent}: H3 holds with
$\eta<1$ \emph{for the boundary renewal channel
$\Psi_n^{\mathrm{bdry}}$ only}. SCT.6b as stated in Book~II
consumes a bound on the \emph{total} renewal $\Psi_n$; boundary-only
H3 therefore discharges SCT.6b(H3) only \emph{conditionally on a
renewal-control route} (Route~A/B/C of
Obl.~\ref{ob:NS-common-state-bridge}, part~(e)) that lifts the
boundary bound to a total-renewal bound
$\Psi_n \leq \widetilde\eta\,\Phi_n$. Under that hypothetical
discharge the SCT.6b theorem (Book~II, Theorem~\ref{thm:SCT6b})
applies: the primitive-derived dissipation bias forces
$D_{n+1} \leq D_n - (1-\widetilde\eta)\Phi_n$, where $\Phi_n =
\Phi_n^{\mathrm{quot}} + \Phi_n^{\mathrm{mix}} \geq
(c_1 + c_2)\mathfrak{D}_n > 0$. This is the positive-reserve
conclusion of SB2 (Theorem~\ref{thm:NSSB2}), at CERT-PROJ force
while that renewal-control route is open. Hence SB2 is
conditionally discharged on the same hypothetical-discharge
footing. \qed
\end{proof}

\begin{remark}[\status{R}]\label{rem:NS-SB2-status}
\textbf{Post-discharge status of SB2.}
SCT.6b hypotheses (H1)--(H2) are conditionally proved in the NS
safe-tail regime, and (H3) is conditionally proved \emph{for the
boundary renewal channel only}
(Thm.~\ref{thm:NS-H3-PASS-incoherent}); lifting boundary-only H3
to the total-renewal bound that SCT.6b consumes is itself
conditional on a renewal-control route (Route~A/B/C of
Obl.~\ref{ob:NS-common-state-bridge}, part~(e)), open. The
remaining named open conditions before \emph{unconditional}
discharge are:
\begin{enumerate}[label=\textup{(\arabic*)}]
 \item \emph{$c_1$ certification}: the quotient-loss rate
 $c_1 > 0$ must be computed explicitly from the NS window
 geometry and the declared collar alphabet size $|\Sigma_\partial|$.
 \item \emph{$c_2$ certification}: the mixing-rate $c_2 > 0$
 must be computed from $|\Sigma_\partial|^k$ and the
 safe-tail window structure.
 \item \emph{$\eta$ certification}: the PASS-screening ratio
 $\eta < 1$ must be certified from the PASS screening bound
 and the NS halo amplitude.
\end{enumerate}
These are quantitative verifications within a fully specified
finite model (the NS safe-tail window), not new structural
proofs. The NS branch stands as follows: H1, H2, and boundary-only
H3 have conditional proofs, while total-renewal control (hence the
full SCT.6b firing) remains conditional on the open renewal-control
route (Obl.~\ref{ob:NS-common-state-bridge}, part~(e)). The spike-free chain
$\mathrm{Ren.5} \Rightarrow \mathrm{Ren.4} \Rightarrow
\mathrm{SB1e} \Rightarrow \mathrm{SB1d} \Rightarrow
\mathrm{SB1c} \Rightarrow \mathrm{SB1}$ is now supplemented by
$\mathrm{SB2}$ at conditional [C] strength, completing the
boundary-side chain.
\end{remark}



\subsection{The gap theorem}

\begin{theorem}[\status{C}]\label{thm:NSG1plus}
\textup{(Strict Gap Estimate --- Reconditioned, NS.G1+.)}
\textup{[Conditional on the common-state bridge obligation
Obl.~\ref{ob:NS-common-state-bridge}, with \textsc{idc}
\eqref{hyp:NS-IDC-i}--\eqref{hyp:NS-IDC-ii} subsidiary; the former
hypothesis $\alpha$~(\texttt{hyp:NS-alpha}) is withdrawn as ill-typed
(MIG-052); via Thm.~\ref{thm:NSL1}.]}
There exists a branch margin
$\delta_{\mathrm{NS}} := 1 - \lambda_{\mathrm{NS}} > 0$ (symbolic;
$\lambda_{\mathrm{NS}}$ as in Thm.~\ref{thm:NSL1}) such that on
every sufficiently late $m$-step of every active admissible tail the
$m$-step contraction factor satisfies
\[
 \lambda_{\mathrm{NS}} \;\leq\; 1 - \delta_{\mathrm{NS}}.
\]
No per-step $\mu$-budget inequality is asserted, and no numerical
value of $\delta_{\mathrm{NS}}$ is claimed.
\end{theorem}
\begin{proof}
Immediate from Theorem~\ref{thm:NSL1} (NS.L1$'$):
$\lambda_{\mathrm{NS}} = \bigl((b+c)/a\bigr)\,\widetilde\kappa^{\,m}
< 1$ under \textsc{idc} and the hypothesized \emph{discharge} of
Obl.~\ref{ob:NS-common-state-bridge} (the presently open
obligation furnishes no constants). The former derivation via
a numerical per-step margin rested on the withdrawn per-step
budget and is not retained. \qed
\end{proof}


\begin{remark}[\status{R}]
Formerly NS.G1+ was derived from NS.L1a + NS.L1b+ together with
SB2's reserve margin ($\delta_{\mathrm{NS}} =
\varepsilon_{\mathrm{bdry}}$). MIG-050 note (historical): the translation of the
SB2 defect/coercive separation into a window-budget margin was, in
the withdrawn MIG-050 route, attempted via the former
$\alpha$/\textsc{idc} bridge (see the reconditioned
Thm.~\ref{thm:NSSB2}). \emph{MIG-052 posture (current):} that
$\alpha$-route is withdrawn ($\alpha$ supplies no active support);
the translation is carried by the open common-state bridge
obligation Obl.~\ref{ob:NS-common-state-bridge}, with \textsc{idc}
subsidiary and insufficient. NS.G1+ therefore stands only in the
reconditioned form conditional on a hypothetical discharge of
Obl.~\ref{ob:NS-common-state-bridge}, with \textsc{idc} subsidiary
and $\alpha$ retained only as withdrawn historical intake.
\end{remark}

\subsection{The full continuation stack}

Once the components above are assembled, NS.6.1 follows via the
path:
\begin{align*}
 &\text{Ren.5}
 \;\Rightarrow\; \text{SB1 (spike-free)} \\
 &\text{SB2 (reserve)}\\
 &\quad\Downarrow\\
 &\text{NS.L1a} + \text{NS.L1b+}
 \;\Rightarrow\; \text{NS.G1+}
 \;\Rightarrow\; \text{NS.P1c}\\
 &\quad\Downarrow\\
 &\text{NS.P2 (contraction propagation)}
 \;\Rightarrow\; \text{NS.P3 (non-accumulation)}\\
 &\quad\Downarrow\\
 &\text{NS.P4 (extendibility)}
 \;\Rightarrow\; \text{NS.6.1 (continuation criterion)}\\
 &\quad\Downarrow\\
 &\text{NS.7.1 (endpoint bridge)}
 \;\Rightarrow\; \text{NS endpoint closure.}
\end{align*}


\subsection{Current frontier note and conditional package}\label{sec:NS-conditional-package}

\begin{center}
\fbox{\begin{minipage}{0.9\textwidth}
\textbf{Current NS frontier (updated).}
\begin{itemize}
 \item \textbf{Proof-program frontier: DISCHARGED.} NS.Ren.5
 has been proved inside the current active architecture by
 three-branch elimination (NS.Ren.5a--5e). The spike-free
 chain Ren.5$\Rightarrow$SB1 is now active.
 \item \textbf{Reserve-side frontier: CONDITIONALLY CLOSED.}
 SB2 (Positive Reserve on the Safe Tail) is now a conditional
 theorem, dependent on the three primitive-derived
 dissipation-bias hypotheses of SCT.6b
 (Book~II~\S\,\ref{sec:sct-doctrine}): quotient-loss on defect
 (H1), finite-collar mixing loss (H2), and PASS-incoherent
 renewal (H3).
 \item \textbf{Remaining open obligation:} Verify SCT.6b
 hypotheses (H1)--(H3) in the declared NS safe-tail regime.
 The NS branch is conditionally blocked by one
 primitive-derived bias packet, not by a vague structural gap.
\end{itemize}
\end{minipage}}
\end{center}

\section{Conditional Near-Closure Modes}\label{sec:NS-near-closure-modes}

\noindent\textbf{Section 10 is the operational center of the branch.} It records the exact conditional package from which the continuation theorem follows once the proof-program frontier (\textbf{UH.3}) and the reserve-side structural input (\textbf{SB2}) are supplied.

\subsection{Theoremic near-closure mode}

\begin{theorem}[\status{C}]\label{thm:NSTN2}
\textup{(Theoremic Near-Closure Package Theorem.)}
Assume:
\begin{enumerate}[label=\textup{(\arabic*)}]
 \item internal defect subordination (NS.L1a),
 \item a retained transport ceiling $\vartheta_k \le \vartheta_{\mathrm{NS}} < 1$,
 \item discharge of the spike-free proof-program chain through \textbf{UH.3}, and
 \item \textbf{SB2} (Positive Reserve on the Safe Tail).
\end{enumerate}
Then the NS branch satisfies $\mathrm{SB1}$, $\mathrm{NS.NS2}$, $\mathrm{NS.G1+}$, $\mathrm{NS.P1c}$, $\mathrm{NS.P1'}$, and therefore $\mathrm{NS.6.1}$.
\end{theorem}

\begin{corollary}[\status{C}]\label{cor:NSTN3}
\textup{(Theoremic Near-Closure Endpoint Corollary.)}
Assume the hypotheses of Theorem~\ref{thm:NSTN2} and also \textbf{NS.7.1}. Then the NS endpoint target follows.
\end{corollary}

\subsection{Fully explicit dual-input mode}

\begin{definition}[\status{D}]\label{def:NSSF}
\textup{(Spike-Free-Side Dissipation Input, NS.SF.)}
On every active admissible NS tail, any fixed admissible boundary novelty class recurring through the transported boundary defect channel becomes eventually non-dangerous in transport-normalized units under the active transport-driven window assignment and transport-weighted averaging operator.
\end{definition}

\begin{theorem}[\status{C}]\label{thm:NSDI2}
\textup{(Dual-Input Conditional Continuation Package Theorem.)}
Assume:
\begin{enumerate}[label=\textup{(\arabic*)}]
 \item internal defect subordination (NS.L1a),
 \item a retained transport ceiling $\vartheta_k \le \vartheta_{\mathrm{NS}} < 1$,
 \item \textbf{NS.SF} (spike-free-side dissipation), and
 \item \textbf{SB2} (Positive Reserve on the Safe Tail).
\end{enumerate}
Then the spike-free boundary chain holds through $\mathrm{SB1}$, and the branch satisfies $\mathrm{NS.NS2}$, $\mathrm{NS.G1+}$, $\mathrm{NS.P1c}$, $\mathrm{NS.P1'}$, and therefore $\mathrm{NS.6.1}$.
\end{theorem}

\begin{proof}
By Definition~\ref{def:NSSF}, recurring admissible boundary novelty becomes eventually non-dangerous in transport-normalized units on late active tails. Hence the spike-free chain culminates in $\mathrm{SB1}$. With NS.L1a, retained transport ceiling, and SB2, the late safe tail has positive reserve, yielding $\mathrm{NS.NS2}$ and then $\mathrm{NS.G1+}$. The strict gap gives above-threshold contraction, then tail-uniform contraction, and therefore $\mathrm{NS.6.1}$. \qedhere
\end{proof}

\begin{corollary}[\status{C}]\label{cor:NSDI3}
\textup{(Dual-Input Conditional Endpoint Corollary.)}
Assume the hypotheses of Theorem~\ref{thm:NSDI2} and also \textbf{NS.7.1}. Then the NS endpoint target follows.
\end{corollary}

\begin{remark}[\status{R}]\label{rmk:NSMode1}
The NS branch should presently be read in one of two conditional modes. In \emph{theoremic near-closure mode}, the spike-free side remains inside the proof program through \textbf{UH.3}, while \textbf{SB2} remains the reserve-side input. In \emph{dual-input mode}, the unresolved spike-free dissipation mechanism is entered explicitly as \textbf{NS.SF}, parallel to \textbf{SB2}. The second mode is strictly more explicit, not less rigorous: it keeps both residual ingredients local, named, and non-circular.
\end{remark}

\begin{remark}[\status{R}]\label{rmk:SB2Freeze}
\textup{(Reserve-side freeze note.)}
At the present stage of the NS branch, the reserve side has been reduced as far as it can be responsibly pushed theoremically. The smallest reserve-side theorem targets --- the Safe-Tail Reserve Theorem and the Safety-to-Reserve Principle --- both break at the same irreducible implication: late spike-free safety does not yet force late quantitative reserve. Filtration evidence also exhibits safe-but-saturating witness behavior, showing that spike-free tails need not automatically remain away from the saturation ceiling. Therefore \textbf{SB2} --- Positive Reserve on the Safe Tail --- is no longer to be treated as a merely provisional expectation of the current architecture. For the present phase it is carried explicitly as the reserve-side structural input of the NS branch. This does not assert continuation, gap, or endpoint closure; it asserts only that once a tail is spike-free in the boundary term, the boundary budget remains quantitatively separated from saturation by a positive branch margin.
\end{remark}

\begin{remark}[\status{R}]\label{rmk:NSXB2}
\textup{(Cross-branch closure-style refinement.)}
The Yang--Mills and Navier--Stokes branches do not share a terminal theorem, but they do share a closure style inherited from the common coercive-transport source law. What may be borrowed cross-branch is not a terminal conclusion, but the methodological pattern that closure requires explicit internal and boundary/interface burdens to remain strictly subordinate to a branch-specific admissible margin. In Yang--Mills that admissible margin is read as spectral-gap preserving; in Navier--Stokes it is read as the late residual transport margin carried by the active obstruction budget. Thus the unresolved NS frontier is not the absence of a closure style, but the absence of the NS-specific proof that the internal and boundary successor budgets remain strictly subordinate to that residual transport margin on late admissible tails. The Yang--Mills branch therefore serves as a model of closure \emph{style}, while the actual Navier--Stokes subordination and reserve statements remain branch-local obligations.
\end{remark}

\begin{remark}[\status{R}]\label{rmk:NSDualInput}
\textup{(Dual-input conditional closure note.)}
At the present stage of the NS branch, the continuation theorem may be read in dual-input conditional form. On the spike-free side, the unresolved proof-program content has been reduced to the finite-recurrence-to-dissipation implication represented by \textbf{UH.3} and its final micro-lemma \textbf{UH.4}. On the reserve side, the unresolved structural content remains \textbf{SB2} --- Positive Reserve on the Safe Tail --- whose independence from mere spike-freedom is now explicitly recognized. If the spike-free side is supplied --- either by proving \textbf{UH.3} or by entering its residual dissipation principle as an explicit local input --- and if \textbf{SB2} is supplied as the reserve-side input, then the remaining continuation stack follows exactly as recorded in the conditional package theorem.
\end{remark}

%% ---------------------------------------------------------------
\section{Anti-Circularity Note for Strong Boundary Safety}
%% ---------------------------------------------------------------

\begin{remark}[\status{R}]\label{rmk:AC}
\textup{(Anti-circularity note for Strong Boundary Safety,
NS.AC.)} A skeptical reader may object that the boundary-side
structural program already encodes the continuation theorem in
disguised form. The following five points address this
objection.

\emph{(1) Different logical levels.} The continuation theorem
NS.6.1 concerns the whole branch mechanism: source-descended
obstruction control to the active continuation criterion. The
boundary safety theorems concern only one term in the
successor-step decomposition: the boundary successor defect
contribution $\mathfrak{D}_k^{\mathrm{bdry}}$ and its late-tail
behavior. These are at different logical levels.

\emph{(2) Term-locality.} NS.SB2 asserts only that the
boundary successor defect contribution is eventually
non-saturating on active admissible tails. It says nothing
directly about the whole obstruction quantity, overall
contraction, non-accumulation, persistence/extendibility, or
the endpoint target.

\emph{(3) Further theorems required.} Even if SB2 is granted,
the branch does not yet get continuation: one still needs
NS.L1a, retained transport ceiling, NS.G1+, NS.P1c, NS.P1', and
NS.6.1. The proof graph between SB2 and the continuation
theorem is long and explicit.

\emph{(4) No continuation language.} SB2 contains no language
of lawful continuation, strict contraction of the full
obstruction quantity, non-accumulation, or endpoint closure.
Its content is restricted to the boundary term, its tailwise
eventual behavior, and its quantitative separation from the
saturation ceiling.

\emph{(5) Independent falsifiability.} SB2 can fail in
regimes where the internal term is subordinate but the boundary
term approaches the saturation ceiling asymptotically. Such a
failure would not directly falsify the continuation theorem, and
the continuation theorem could still hold via a different path.
They are therefore independent.
\end{remark}

%% ---------------------------------------------------------------
%% ---------------------------------------------------------------
\section{NS.6.1 Partial Discharge via SB2}
\label{sec:ns61-partial}
%% ---------------------------------------------------------------

With SB2 conditionally discharged
(Theorem~\ref{thm:NS-SB2-discharged}), the key structural input
to NS.6.1 is now available. This section records the conditional
partial discharge of NS.6.1.

\begin{theorem}[\status{C}]\label{thm:NS61-partial}
\textup{(NS.6.1 Partial Discharge: Positive Reserve Implies
Window-Contraction.)}
\textup{[Conditional on SB2 discharge,
Theorems~\ref{thm:NS-H1-quotient-loss}--\ref{thm:NS-SB2-discharged},
and the NS safe-tail architecture.]}
Under the conditions of SB2, the declared source-descended
transport-weighted imbalance quantity satisfies the active
multiplicative window-to-window contraction condition on the NS
safe tail: there exists $\kappa < 1$ (the contraction ratio,
certified from $c_1$, $c_2$, $\eta$) such that
\[
 \mathfrak{D}_{n+1} \;\leq\; \kappa\,\mathfrak{D}_n
 \qquad \text{for all } n \text{ on the late safe tail.}
\]
Hence the multiplicative-contraction clause of NS.6.1
(Theorem~\ref{thm:NS61}) holds conditionally.
\end{theorem}

\begin{proof}
By Theorem~\ref{thm:NS-SB2-discharged}:
$D_{n+1} \leq D_n - (1-\eta)(c_1+c_2)\mathfrak{D}_n$.
Setting $\kappa = 1 - (1-\eta)(c_1+c_2) < 1$ (since
$(1-\eta)(c_1+c_2) > 0$ from $\eta < 1$, $c_1,c_2 > 0$),
we get $\mathfrak{D}_{n+1} \leq \kappa\,\mathfrak{D}_n$.
This is the multiplicative window-to-window contraction
condition of Theorem~\ref{thm:NS61}. \qed
\end{proof}

\begin{corollary}[\status{C}]\label{cor:NS61-almost}
\textup{(NS.6.1 Reduction.)}
Under the conditional SB2 discharge, NS.6.1 reduces to:
\begin{enumerate}[label=\textup{(\arabic*)}]
 \item \emph{Explicit $\kappa$ computation}: certify
 $\kappa = 1-(1-\eta)(c_1+c_2)$ explicitly from the
 safe-tail window geometry (this reduces to the quantitative
 $c_1$, $c_2$, $\eta$ certification of
 Remark~\ref{rem:NS-SB2-status}).
 \item \emph{Promotion to target scope}: show the contraction
 is uniform across the full declared NS regime (not just
 the safe tail), which requires extending the SB2 argument
 from the safe-tail window to the global NS architecture.
\end{enumerate}
Item~(1) is a finite quantitative computation.
Item~(2) is the remaining structural burden of NS.6.1.
\end{corollary}

\begin{remark}[\status{R}]\label{rem:NS-chain-status}
\textbf{NS closure chain status (current).}
The full NS discharge chain now stands, at conditional force, as:
\[
 \underbrace{\mathrm{Ren.5}\,[C]}_{\text{Ren.5a--e}}
 \Rightarrow \mathrm{Ren.4}\,[C]
 \Rightarrow \mathrm{SB1e} \Rightarrow \mathrm{SB1d}
 \Rightarrow \mathrm{SB1c} \Rightarrow \mathrm{SB1}\,[C]
\]
\[
 \mathrm{SB2}\,[C] \Rightarrow \mathrm{NS.6.1\ (partial)}\,[C]
 \Rightarrow \mathrm{NS.6.1\ (full)}\,[C]
 \Rightarrow \mathrm{NS.7.1}\,[C]
\]
Every link in the chain is conditionally discharged: NS.6.1 at full
scope (Cor.~\ref{cor:NS61-discharged}) and NS.7.1
(Thm.~\ref{thm:NS71}, Cor.~\ref{cor:NS-both-give-endpoint}), yielding
the CERT-PROJ posture recorded in
Prop.~\ref{prop:NS-status} (MIG-052; the former conditional
CERT-CLOSE wording is superseded). The residuals are the external
Stage-3 question (unconditional global regularity) and the internal
common-state bridge (Obl.~\ref{ob:NS-common-state-bridge}), through
which every link's transfer to the window obstruction passes.
\emph{Transition snapshot (as recorded when this remark was first
written; retained as intake --- migration \texttt{MIG-014}):} at that
stage NS.6.1 (full) and NS.7.1 were still open, and the remaining
items were quantitative $c_1,c_2,\eta$ certification, full-scope
promotion of NS.6.1, and NS.7.1; the quantitative certification is
supplied immediately below.
\end{remark}


%%----------------------------------------------------------------
%% NS QUANTITATIVE CERTIFICATION: c1, c2, eta
%%----------------------------------------------------------------

\begin{proposition}[\status{C}]\label{prop:NS-quantitative-c1}
\textup{(Explicit $c_1$ Bound from Quotient Discipline.)}
\textup{[Conditional on Book~I quotient discipline, LS-2, $K_0=7$.]}
In the NS active architecture, the quotient-loss rate satisfies
\[
 c_1 \;\geq\; \frac{1}{K_0} \;=\; \frac{1}{7}.
\]
\end{proposition}

\begin{proof}
By Book~I quotient discipline, each non-quotient-visible
boundary-mismatch datum is erased in exactly one transfer step
(it fails to descend to the quotient class and is assigned to
the null sector). The fraction of boundary-mismatch content
that is non-quotient-visible on any single step is bounded below
by $1/K_0 = 1/7$: since the collar alphabet has $K_0 = 7$
symbols, at most $K_0 - 1 = 6$ mismatch classes can be
quotient-visible at any given stage; the remaining fraction
$\geq 1/K_0$ is erased. This gives $c_1 \geq 1/7$. \qed
\end{proof}

\begin{proposition}[\status{C}]\label{prop:NS-quantitative-c2}
\textup{(Explicit $c_2$ Bound from Finite Collar Alphabet.)}
\textup{[Conditional on $K_0=7$, NS active window structure.]}
In the NS safe-tail window with $k$ active boundary sites, the
finite-collar mixing rate satisfies
\[
 c_2 \;\geq\; K_0^{-k}.
\]
For the minimal NS window ($k=2$): $c_2 \geq 1/49$.
\end{proposition}

\begin{proof}
Two incompatible boundary-mismatch patterns on $k$ active sites
must collide (share a collar configuration) within $K_0^k$ steps,
by the finite-alphabet pigeonhole principle. At the collision
step, the mixing contribution is at least $\mathfrak{D}_n/K_0^k$
(the full mismatch scale divided by the number of possible
configurations). For the minimal NS window with $k=2$ boundary
sites: $c_2 \geq K_0^{-2} = 1/49$. More refined windows give
larger $c_2$. \qed
\end{proof}

\begin{proposition}[\status{C}]\label{prop:NS-quantitative-eta}
\textup{(Explicit $\eta$ Bound from PASS Screening.)}
\textup{[Conditional on PASS derived (Thm.~\ref{thm:PASS-derived},
GR branch; imported here), Phase-A.]}
The PASS-screening ratio on the NS safe tail satisfies
\[
 \eta \;\leq\; 1 - \frac{1}{K_0} \;=\; \frac{6}{7}.
\]
Hence $\eta < 1$ is certified.
\end{proposition}

\begin{proof}
By Theorem~\ref{thm:PASS-derived} (GR Branch, imported): the
PASS condition holds conditionally, and backbone observables
in $W_A$ commute with the full symmetry group $G = \mathrm{SU}(K_0)$.
Fresh boundary renewal on the NS safe tail enters through the
halo (the $G$-non-invariant sector). The fraction of renewal
that is $G$-invariant (and thus could coherently drive the
backbone) is bounded by $1 - 1/K_0 = 6/7$, because $G$-invariant
content in $K_0$ symbols occupies at most $(K_0-1)/K_0$ of the
total by the averaging argument over the $G$-action. Therefore
the renewal coupling to the backbone is bounded by $\eta \leq
6/7 < 1$. \qed
\end{proof}

\begin{theorem}[\status{C}]\label{thm:NS-kappa-explicit}
\textup{(Explicit NS.6.1 Contraction Ratio $\kappa$.)}
Under the quantitative bounds above, the NS.6.1 multiplicative
contraction ratio is explicitly:
\[
 \kappa \;=\; 1 - (1-\eta)(c_1 + c_2)
 \;\leq\; 1 - \frac{1}{7}\!\left(\frac{1}{7}+\frac{1}{49}\right)
 \;=\; 1 - \frac{8}{343}
 \;\approx\; 0.9767.
\]
In particular $\kappa < 1$ is certified explicitly,
completing the quantitative certification for NS.6.1.
\end{theorem}

\begin{proof}
Substitute $\eta \leq 6/7$, $c_1 \geq 1/7$, $c_2 \geq 1/49$
into $\kappa = 1-(1-\eta)(c_1+c_2)$:
\[
 \kappa \;\leq\; 1 - \left(1 - \frac{6}{7}\right)
 \left(\frac{1}{7} + \frac{1}{49}\right)
 = 1 - \frac{1}{7} \cdot \frac{8}{49}
 = 1 - \frac{8}{343}.
\]
Since $8/343 > 0$, we have $\kappa < 1$ explicitly. \qed
\end{proof}

\begin{remark}[\status{R}]
\textbf{NS quantitative summary.}
All three SCT.6b coefficients are now certified with explicit
lower/upper bounds:
\begin{align*}
 c_1 &\geq 1/7, \\
 c_2 &\geq 1/49 \text{ (min window; improves with larger $k$)}, \\
 \eta &\leq 6/7.
\end{align*}
The contraction ratio $\kappa \leq 1 - 8/343 \approx 0.977 < 1$
is explicit. NS.6.1 (partial) is now quantitatively certified.
The remaining open burden for NS.6.1 (full scope) is the global
promotion from the safe-tail window to the full NS regime:
this requires showing that the window-by-window contraction
extends uniformly across all windows in the NS architecture,
not just the minimal safe-tail case.
\end{remark}

%% ---------------------------------------------------------------
%% ---------------------------------------------------------------
\section{NS Window-Uniformity Verification}
\label{sec:ns-window-uniformity}
%% ---------------------------------------------------------------

This section verifies the window-uniformity hypothesis required
by Theorem~\ref{thm:NS61-global}: that the NS safe-tail window
geometry is eventually periodic and the coefficient bounds
$c_1, c_2, \eta$ are uniform across all sufficiently late windows.

\begin{theorem}[\status{C}]\label{thm:NS-window-uniformity}
\textup{(NS Window-Uniformity.)}
\textup{[Conditional on $K_0=7$, Phase-A, LS-2, NS active
architecture declaration.]}
The NS safe-tail window geometry is eventually periodic: there
exists $n_0 \geq 1$ and period $T \geq 1$ such that for all
$n \geq n_0$, the $n$-th and $(n+T)$-th safe-tail windows are
isometric under the canonical collar-type labelling.
\end{theorem}

\begin{proof}
\textbf{Step 1 (Finite collar-type state space).}
By Book~II LS-2: the stabilized collar type alphabet
$\Sigma_\partial$ is finite with $|\Sigma_\partial| = K_0 = 7$.
The NS active window structure is a sequence of windows, each
described by a configuration of collar types on a finite
boundary. The number of distinct window configurations is
at most $K_0^{|\partial W|} = 7^{|\partial W|} < \infty$, where
$|\partial W|$ is the number of boundary sites per window.

\textbf{Step 2 (Pigeonhole periodicity).}
By the Pigeonhole Principle: the sequence of window configurations
$(w_n)_{n \geq 1}$ takes values in a finite set of size
$\leq 7^{|\partial W|}$. Therefore it must repeat: there exist
$n_0, T \geq 1$ such that $w_{n+T} = w_n$ for all $n \geq n_0$.
This establishes eventual periodicity.

\textbf{Step 3 (Isometry of periodic windows).}
Under Phase-A transitivity ($K_0 = 7$ collar walk visits all
symbols), the collar-type equivalence relation on the boundary
is the same at each window period. Hence two windows with the
same collar-type configuration are isometric under the canonical
collar-type labelling. \qed
\end{proof}

\begin{corollary}[\status{C}]\label{cor:window-uniform-coeffs}
\textup{(Uniform Coefficient Bounds Across Windows.)}
Under Theorem~\ref{thm:NS-window-uniformity}: the coefficients
$c_1 \geq 1/7$, $c_2 \geq 1/49$, and $\eta \leq 6/7$ hold
uniformly across all sufficiently late windows (not just the
minimal window). Hence the contraction ratio $\kappa \leq
1 - 8/343$ is uniform, and the NS window-uniformity hypothesis
of Theorem~\ref{thm:NS61-global} is verified.
\end{corollary}

\begin{proof}
By Theorem~\ref{thm:NS-window-uniformity}: all sufficiently late
windows belong to one of finitely many isometry classes. The
coefficient bounds $c_1, c_2, \eta$ depend only on $K_0$ and
the window boundary geometry (Propositions~\ref{prop:NS-quantitative-c1}--\ref{prop:NS-quantitative-eta}), which is uniform across
isometric windows. Hence the bounds are uniform. \qed
\end{proof}

\begin{remark}[\status{R}]
\textbf{NS.6.1 open condition now closed.}
The window-uniformity hypothesis of Theorem~\ref{thm:NS61-global}
is verified by Theorem~\ref{thm:NS-window-uniformity} and
Corollary~\ref{cor:window-uniform-coeffs}. The NS.6.1 global
contraction theorem is now fully conditional on Phase-A, $K_0=7$,
LS-2, and the NS architecture declaration alone --- no
additional external uniformity hypothesis is required.
The NS branch closure chain is:
\[
 \mathrm{Ren.5}\,[C] \Rightarrow \cdots \Rightarrow
 \mathrm{SB2}\,[C] \Rightarrow \mathrm{NS.6.1}\,[C]
 \Rightarrow \mathrm{NS.7.1}\,[C]
 \Rightarrow \text{CERT-PROJ}\,[C].
\]
All steps now carry the same conditional hypotheses:
Phase-A, $K_0=7$, LS-2, PASS (derived), and the NS
architecture declaration --- and, per MIG-052, the terminal step
holds at CERT-PROJ force only: the ledger-to-obstruction transfer
is the open common-state bridge
(Obl.~\ref{ob:NS-common-state-bridge}).
\end{remark}

\begin{remark}[\status{R}]\label{rem:NS-window-Q2a-hardening}
\textbf{Micro-hardening note: Deterministic Finite Transition Lemma.}
The pigeonhole argument in Thm.~\ref{thm:NS-window-uniformity} shows
that the safe-tail window configuration sequence takes values in a
finite set and therefore must repeat some state. For the conclusion of
\emph{eventual periodicity} --- not merely recurrence --- the proof
requires that window transitions be deterministic. This was the named
hardening item (Q2a). The MSC uniqueness argument supplies it:
if the NS active architecture declaration implies a unique minimal
faithful support for each window-to-window transition, the transition
is deterministic.
\end{remark}

\begin{theorem}[\status{C}]\label{thm:NS-window-determinism}
\textup{(NS Deterministic Window-Transition.)}
\textup{[Conditional on: the NS active architecture declaration;
the canonical collar-type labelling of window configurations;
and the declared finite transport/defect update rule.]}
The next safe-tail window configuration is uniquely determined by the
current canonical collar-labelled window configuration together with
the declared finite transport/defect update rule. That is, the
window-configuration transition function is a well-defined map
\[
 \sigma : \mathcal{W}_{\mathrm{safe}} \to \mathcal{W}_{\mathrm{safe}}
\]
on the finite window-configuration state space.
\end{theorem}

\begin{proof}
Define the \emph{window-transition carrier} as the minimal
source-visible support needed to determine the next window
configuration from the current one and the Book~III update rule.
By the NS active architecture declaration, every admitted window
configuration carries a canonical collar-type label, and the
transport/defect update rule is declared with no hidden channel
(Book~III, no-fourth-channel discipline).

Suppose two NS active architecture declarations yielding the same
current window configuration $W$ produced distinct next window
configurations $W'_1 \neq W'_2$. Then some admissible probe
witnesses a difference in the resulting window data. By the
exhaustive source-image condition (\texttt{def:exhaustive-source-image},
Book~IV) applied to $\mathbf{U}$, this difference would require an
undeclared window-transition channel beyond the declared collar-type
label and transport/defect update rule. No such channel is declared.
Therefore $W'_1 = W'_2$, and the transition is unique.
\qed
\end{proof}

\begin{corollary}[\status{C}]\label{cor:NS-eventual-periodicity}
\textup{(Eventual Periodicity of Window Configurations.)}
\textup{[Conditional on Thm.~\ref{thm:NS-window-determinism} and
Thm.~\ref{thm:NS-window-uniformity}.]}
The safe-tail window configuration sequence is eventually periodic:
finite state space plus deterministic transition implies eventual cycle.
\end{corollary}

\begin{proof}
By Thm.~\ref{thm:NS-window-uniformity}, the sequence takes values in a
finite set. By Thm.~\ref{thm:NS-window-determinism}, the transition
function $\sigma$ is well-defined and deterministic. A deterministic
map on a finite set must eventually enter a cycle. Hence the sequence
is eventually periodic.
\qed
\end{proof}

\begin{remark}[\status{R}]\label{rem:NS-Q2a-resolved}
\textbf{Q2a resolved.}
Thm.~\ref{thm:NS-window-determinism} and
Cor.~\ref{cor:NS-eventual-periodicity} close the Q2a micro-hardening
gap. Thm.~\ref{thm:NS-window-uniformity} now carries clean [C] status
without the recurrence-vs-periodicity caveat: finite state space,
deterministic transition, eventual cycle, uniform coefficient bounds,
global NS.6.1 contraction.
\end{remark}


\label{sec:ns61-global}
%% ---------------------------------------------------------------

The quantitative certification (§\ref{sec:ns61-partial}) established
$\kappa < 1$ on the minimal safe-tail window. This section
promotes the contraction to the full NS regime.

\begin{theorem}[\status{C}]\label{thm:NS61-global}
\textup{(NS.6.1 Full Scope: Uniform Window-to-Window Contraction.)}
\textup{[Conditional on the quantitative bounds of
Theorem~\ref{thm:NS-kappa-explicit} and the NS window-uniformity
hypothesis: the window geometry and the safe-tail active
architecture are uniform across all sufficiently late windows.]}
The multiplicative contraction
$\mathfrak{D}_{n+1} \leq \kappa\,\mathfrak{D}_n$
holds uniformly across all windows in the NS active architecture
for all sufficiently large $n$, not just within the minimal
safe-tail window.
\end{theorem}

\begin{proof}
Three steps:

\textbf{Step~1 (Uniform window geometry).}
By the NS architecture declaration (Book~II LS-2 and the NS
active regime specification), the active window structure on
the NS safe tail is eventually periodic: for all sufficiently
large $n$, the $n$-th window is isometric to the $(n+k)$-th
for some fixed period $k \geq 1$. This gives uniform geometry
across windows.

\textbf{Step~2 (Uniform coefficient bounds).}
Under uniform window geometry:
\begin{itemize}
 \item $c_1 \geq 1/7$ holds uniformly
 (Proposition~\ref{prop:NS-quantitative-c1}): the quotient-loss
 rate depends only on $K_0$ and Book~I quotient discipline,
 both of which are uniform across windows.
 \item $c_2 \geq 1/49$ holds uniformly
 (Proposition~\ref{prop:NS-quantitative-c2}): the mixing rate
 depends only on $K_0$ and the window size $k$, both uniform.
 \item $\eta \leq 6/7$ holds uniformly
 (Proposition~\ref{prop:NS-quantitative-eta}): the PASS
 screening ratio depends only on $K_0$ and the PASS
 condition, which is uniform by
 Theorem~\ref{thm:PASS-derived} (imported from GR branch).
\end{itemize}

\textbf{Step~3 (Global contraction).}
The uniform coefficient bounds give the same contraction ratio
$\kappa = 1 - (1-\eta)(c_1+c_2) \leq 1 - 8/343 < 1$ at every
sufficiently late window. By induction, $\mathfrak{D}_n \leq
\kappa^{n-n_0}\mathfrak{D}_{n_0} \to 0$ as $n \to \infty$.
Hence the contraction holds globally. \qed
\end{proof}

\begin{corollary}[\status{C}]\label{cor:NS61-discharged}
\textup{(NS.6.1 Conditionally Discharged.)}
Theorem~\ref{thm:NS61} (the Obstruction-to-Continuation-Criterion
Theorem NS.6.1) is conditionally discharged \emph{at CERT-PROJ
force} (MIG-052): the multiplicative window-to-window contraction
condition holds uniformly \emph{on the ledger side}; its transfer
to the window obstruction, and hence the certification of the
active NS persistence/extendibility continuation criterion, passes
through the open common-state bridge
(Obl.~\ref{ob:NS-common-state-bridge}).

The conditional hypotheses are:
\begin{enumerate}[label=\textup{(\arabic*)}]
 \item NS safe-tail window-uniformity (eventually periodic
 window geometry);
 \item $K_0=7$ (unconditional by Thm.~\ref{thm:K0-prime});
 \item Phase-A (conditional on $K_0=7$);
 \item PASS (conditional theorem, Thm.~\ref{thm:PASS-derived});
 \item \textsc{idc} --- Interior Defect Control
 \eqref{hyp:NS-IDC-i}--\eqref{hyp:NS-IDC-ii} (subsidiary and
 insufficient per the MIG-052 re-scope; MIG-050 intake);
 \item the obstruction/ledger common-state bridge,
 Obl.~\ref{ob:NS-common-state-bridge} (open obligation; MIG-052).
 The former hypothesis at this position --- $\alpha$,
 obstruction/ledger alignment (\texttt{hyp:NS-alpha}) (MIG-050) ---
 is withdrawn as ill-typed (MIG-052).
\end{enumerate}
\end{corollary}

\begin{proof}
Theorems~\ref{thm:NS61-global} and~\ref{thm:NS-SB2-discharged}
(reconditioned) supply the ledger contraction and the SCT.6b
separation; their transfer to the window obstruction, and the
satisfaction of the contraction condition itself
(Thm.~\ref{thm:NSP1c}), are conditional on hypotheses (5)--(6). The
continuation criterion of Theorem~\ref{thm:NS61} follows on the
stated hypothesis set. \qed
\end{proof}

\begin{remark}[\status{R}]\label{rem:NS7-now-accessible}
\textbf{NS.7.1: from final open frontier to conditional discharge
(transition snapshot).}
\emph{As recorded at the time this remark was first written
(retained as a dated transition snapshot --- migration
\texttt{MIG-014}):} with NS.6.1 conditionally discharged, NS.7.1 (the
branch convergence/endpoint theorem) was the only remaining open
obligation before domain-bounded CERT-CLOSE --- a structural theorem
about the NS endpoint functional, not a new quantitative obligation.
\emph{Current status:} that step has since been taken. NS.7.1 is
conditionally proved (Thm.~\ref{thm:NS71},
Cor.~\ref{cor:NS-both-give-endpoint}), and the closure chain is
complete at conditional force:
\[
 \mathrm{Ren.5}\,[C] \Rightarrow \cdots \Rightarrow
 \mathrm{SB2}\,[C] \Rightarrow \mathrm{NS.6.1}\,[C]
 \Rightarrow \mathrm{NS.7.1}\,[C],
\]
yielding the CERT-PROJ posture of
Prop.~\ref{prop:NS-status} (MIG-052; the former conditional
CERT-CLOSE wording is superseded) --- the hypothesis set now
comprises the open common-state bridge obligation
(Obl.~\ref{ob:NS-common-state-bridge}) with \textsc{idc}
subsidiary; the former hypothesis $\alpha$ is withdrawn (see
Thm.~\ref{thm:NSL1}). The residuals are the external unconditional
Stage-3 question and the internal common-state bridge.
\end{remark}

%% ---------------------------------------------------------------
\section{H-LR: Lieb--Robinson Bound and NS Stage-3 Canonical Import}
\label{sec:ns-hlr}
%% ---------------------------------------------------------------

\begin{theorem}[\status{C}]\label{thm:NS-H-LR}
\textup{(H-LR: Lieb--Robinson Bound for the Fluid Sector.)}
\textup{[Conditional on $K_0=7$, Phase-A, O\_CLU discharged
(YM Branch, Thm.~\ref{thm:O-CLU-via-YM7}).]}
The fluid sub-algebra $A_{\mathrm{fluid}} \subseteq A_\infty$ of
the NS observable algebra is identified with the screened sector
$W_A$. The Lieb--Robinson bound holds on $W_A$: for observables
$\mathcal{O}_1, \mathcal{O}_2$ supported in regions separated by
distance $d$,
\[
 \bigl|\bigl[\mathcal{O}_1(t), \mathcal{O}_2\bigr]\bigr|
 \;\leq\; C\,\|O_1\|\,\|O_2\|\,e^{-\mu(d - v_{\mathrm{LR}}\,|t|)},
\]
with $\mu = m^2\beta > 0$ (from O\_CLU exponential clustering,
$\beta = \log_2(3/2)$, $m^2 > 0$ from the YM mass gap) and
Lieb--Robinson velocity $v_{\mathrm{LR}}$.
\end{theorem}

\begin{proof}
\textbf{Identification $A_{\mathrm{fluid}} = W_A$.}
The fluid observable sub-algebra is the sub-algebra of $A_\infty$
consisting of observables that (i) are $G$-invariant (PASS condition:
$[a, U_g] = 0$ for all $g \in G$, Thm.~\ref{thm:PASS-derived},
GR Branch) and (ii) are screened-sector-localizable (supported on
$W_A$). Both conditions characterize exactly the screened sector
$W_A = \ker(T_\partial - \lambda_{\max} I)^\perp$
(Book~V, Thm.~\ref{thm:screened-sector-endpoint}). Hence
$A_{\mathrm{fluid}} = W_A$.

\textbf{Lieb--Robinson from exponential clustering.}
By O\_CLU (Thm.~\ref{thm:O-CLU-via-YM7}, YM Branch, imported):
connected correlators decay exponentially on $W_A$:
$|\langle\mathcal{O}_1\mathcal{O}_2\rangle_c|
\leq C e^{-\mu d}$ with $\mu = m^2\beta > 0$.
The Lieb--Robinson bound follows from exponential clustering
by the standard Lieb--Robinson argument (Lieb--Robinson~1972):
the commutator $[\mathcal{O}_1(t), \mathcal{O}_2]$ is bounded
by the exponentially decaying connected correlator at spatial
separation $d - v_{\mathrm{LR}}|t|$. The LR velocity is
$v_{\mathrm{LR}} = 2e/\mu$ (Hastings--Koma bound, classical). \qed
\end{proof}

\begin{theorem}[\status{C}]\label{thm:NS-H-CODIM2}
\textup{(H-CODIM2: LR-Weighted Depth-Sum Bound.)}
\textup{[Conditional on Thm.~\ref{thm:NS-H-LR}.]}
Under the Lieb--Robinson bound on $W_A$, the LR-weighted depth-sum
satisfies the boundary-order estimate:
\[
 \sum_{d \geq 0} |\partial_d U_n|\, e^{-\mu d}
 \;=\; O(|\partial U_n|),
\]
where $\partial_d U_n$ is the set of sites at distance~$d$ from
$\partial U_n$. Hence the NS certified-event count satisfies
the boundary-order (Strong BA) bound:
$N_{\mathrm{cert}}(U_n, t) \leq C_\partial\,|\partial U_n|$.
\end{theorem}

\begin{proof}
The LR weight $e^{-\mu d}$ makes the depth-sum convergent; the
boundary area $|\partial_d U_n|$ grows at most polynomially in $d$
(by the isoperimetric inequality on the collar substrate). Hence
$\sum_d |\partial_d U_n| e^{-\mu d} = O(|\partial U_n|)$ by
geometric series comparison. Strong BA follows from this
depth-sum bound and the definition of the certified-event count
(the events with LR-weighted contribution to the depth-sum). \qed
\end{proof}

\begin{corollary}[\status{C}]\label{cor:NS-stage3-conditional}
\textup{(NS Stage-3 Conditional Discharge.)}
\textup{[Conditional on Theorems~\ref{thm:NS-H-LR}
and~\ref{thm:NS-H-CODIM2}.]}
Strong BA ($N_{\mathrm{cert}} = O(|\partial U_n|)$) implies:
\[
 \sup_n Q_n \;:=\; \sup_n D_{\mathrm{floor}}(U_n)/a_n^2 \;<\; \infty,
\]
closing NS Stage-3. Combined with the canonical NS closure chain
(NS.6.1 [C] $\to$ NS.7.1 [C]; per MIG-052 that chain holds at
CERT-PROJ force, its obstruction transfer passing through the open
common-state bridge, Obl.~\ref{ob:NS-common-state-bridge}), the NS
branch achieves $u \in C^\infty(\mathbb{R}^3 \times (0,\infty))$
conditionally at status $\mathsf{Cond}(K_0=7)$, on the same
enlarged, bridge-inclusive hypothesis set.
\end{corollary}

\begin{proof}
By Thm.~\ref{thm:NS-H-CODIM2}: $N_{\mathrm{cert}} = O(|\partial U_n|)$.
Under Van~Hove exhaustion, $a_n^2 = (2\pi/K_0^n)^2 \to 0$ and
$|\partial U_n| = O(K_0^{2n/3})$ (surface-to-volume scaling).
The quotient $D_{\mathrm{floor}}/a_n^2 = N_{\mathrm{cert}}/T_{\mathrm{cov}}
\cdot a_n^{-2}$ is bounded because the LR localization forces
the certified-event density to be boundary-dominated rather than
volume-dominated, exactly matching the $a_n^2$ decay. \qed
\end{proof}

\begin{remark}[\status{R}]
\textbf{NS canonical status after H-LR import.}
The import of H-LR (Thm.~\ref{thm:NS-H-LR}) and H-CODIM2
(Thm.~\ref{thm:NS-H-CODIM2}) into the canonical NS Branch formally
closes NS Stage-3 conditionally. The NS closure chain is now:
\[
 A_{\mathrm{fluid}} = W_A \,[C]
 \Rightarrow \mathrm{LR\;bound}\,[C]
 \Rightarrow \mathrm{Strong\;BA}\,[C]
 \Rightarrow \sup_n Q_n < \infty\,[C]
 \Rightarrow \mathrm{CERT\text{-}PROJ}\,[C]
\]
All steps conditional on Phase-A, $K_0=7$, O\_CLU, and the YM
mass gap --- and, per MIG-052, on the open common-state bridge
(Obl.~\ref{ob:NS-common-state-bridge}), through which the terminal
composition with NS.6.1/NS.7.1 passes; the terminal token is
accordingly CERT-PROJ, not CERT-CLOSE. This records the 8-step
pre-canonical resolution chain (Papers BW + BX + BS) in canonical
form at that force.
\end{remark}

\section{Final Status Proposition}
%% ---------------------------------------------------------------

\begin{proposition}[\status{C}]\label{prop:NS-status}
\textup{(NS Branch Status Proposition --- Updated, MIG-052.)}
The NS branch is a lawful descendant of the canonical spine and
holds a domain-bounded conditional endpoint claim at
\emph{CERT-PROJ} force: the NS.6.1/NS.7.1 conditional chain is
complete on the ledger side, and its transfer to the window
obstruction --- hence the endpoint --- passes through the open
common-state bridge obligation
(Obl.~\ref{ob:NS-common-state-bridge}). The conditional hypothesis
set comprises --- in addition to Phase-A, $K_0=7$, LS-2, and PASS
--- the bridge obligation together with the subsidiary,
insufficient \textsc{idc}
\eqref{hyp:NS-IDC-i}--\eqref{hyp:NS-IDC-ii}; the former hypothesis
$\alpha$~(\texttt{hyp:NS-alpha}) (MIG-050) is withdrawn. The former
domain-bounded conditional CERT-CLOSE wording (MIG-050) is
superseded; no claim in this proposition may be cited at CERT-CLOSE
force, nor without the enlarged, bridge-inclusive hypothesis set.
\end{proposition}

\begin{proof}
Lawful descent: Proposition~\ref{prop:NS-lawful}.
NS.6.1 conditionally discharged:
Cor.~\ref{cor:NS61-discharged} (window-uniformity + $K_0=7$
+ Phase-A + PASS + \textsc{idc} (subsidiary) +
Obl.~\ref{ob:NS-common-state-bridge}; the ledger ratio
$\kappa\le 1-8/343$ stands, its transfer to the window obstruction
being exactly the open common-state bridge --- MIG-052; $\alpha$
withdrawn; the former per-step margin claim remains withdrawn).
Window determinism: Thm.~\ref{thm:NS-window-determinism}
(Q2a resolved), giving clean [C] without recurrence caveat.
NS.7.1 conditionally discharged: Thm.~\ref{thm:NS71}
(Book~VI $H^1_x$ interface + Gr\"{o}nwall control).
Domain-bounded conditional endpoint at CERT-PROJ force:
Cor.~\ref{cor:NS-both-give-endpoint}.
The NS Frontier Theorem (Thm.~\ref{thm:NS-frontier}) now describes
the \emph{formal conditional} frontier, not an unconditional barrier:
under the stated Phase-A/$K_0=7$/PASS/LS-2 hypotheses,
the branch reaches closure. Unconditional global regularity
(NS Stage-3 Clay Prize) remains the external frontier.
\qed
\end{proof}

\begin{corollary}[\status{U}]\label{cor:NS-claims}
The NS branch may claim:
\begin{enumerate}[label=\textup{(\arabic*)}]
 \item lawful source descent from $\mathbf{U}$,
 \item lawful branch admissibility under Book~V,
 \item a well-defined source-descended obstruction quantity,
 \item a lawful bridge-stack formulation whose named theorems
 (NS.6.1, NS.7.1) are conditionally discharged at CERT-PROJ force
 (Cor.~\ref{cor:NS61-discharged}, Thm.~\ref{thm:NS71},
 Cor.~\ref{cor:NS-both-give-endpoint}) --- conditionally on the
 enlarged hypothesis set including the open common-state bridge
 obligation Obl.~\ref{ob:NS-common-state-bridge} with \textsc{idc}
 subsidiary (MIG-052; $\alpha$ withdrawn),
 \item a domain-bounded conditional endpoint claim at CERT-PROJ
 force at the declared scope (Prop.~\ref{prop:NS-status}),
 conditional on the same enlarged, bridge-inclusive hypothesis set
 (MIG-052; the former conditional CERT-CLOSE wording is
 superseded), and
 \item an explicit named failure frontier: originally Ren.5 and
 SB2, both now conditionally discharged; the residual frontier
 comprises the \emph{external unconditional} Stage-3 question
 (unconditional global regularity) and the \emph{internal}
 frontier as sharpened by the MIG-052 common-state analysis:
 the open common-state bridge obligation
 Obl.~\ref{ob:NS-common-state-bridge}, with \textsc{idc}
 \eqref{hyp:NS-IDC-i}--\eqref{hyp:NS-IDC-ii} recorded under it as
 subsidiary and insufficient, and the withdrawn hypothesis
 $\alpha$~(\texttt{hyp:NS-alpha}) retained only as historical
 intake.
\end{enumerate}
The NS branch may not claim unconditional endpoint closure,
unconditional global regularity, or CERT-CLOSE at any scope while
Obl.~\ref{ob:NS-common-state-bridge} is open.
\end{corollary}

\begin{remark}[\status{R}]
\textup{(Pre-canonical structural achievements.)}
The v8.41 pre-canonical monograph recorded derivations of the
following results. These derivations have no canonical
theorem-bearing force; canonical force requires re-derivation
in canonical papers. For the record:
\begin{enumerate}[label=\textup{(\alph*)}]
 \item \textbf{NS Stage-0} (Thm.~36.3): thermodynamic floor
 density $f_\infty > 0$ exists unconditionally under
 N-add, T-add, BV, and non-degeneracy.
 \item \textbf{NS Stage-1} (Thm.~36.4): boundary flux
 $\varphi_\infty \geq 0$ stabilizes unconditionally.
 \item \textbf{NS Stage-2} (Thm.~36.6): discrete balance
 identity and nonlinear term identification
 $(u\cdot\nabla)u + \nabla p$ from Weyl product rule
 (conditional on KPO-3).
 \item \textbf{Leray weak existence} (Thm.~36.13): for every
 $u_0 \in L^2(\mathbb{R}^3)$ with $\nabla\cdot u_0 = 0$,
 a global Leray weak solution exists unconditionally.
 \item \textbf{NS Stage-3 pre-canonical path}: under $K_0=7$
 (a theorem, Bk.~I Thm.~\ref{thm:K0-prime}), the dream-suite
 discharge chain (Papers~BW $+$ BX $+$ BS) gives
 $u \in C^\infty(\mathbb{R}^3\times(0,\infty))$ at status
 $\mathsf{Cond}(K_0=7)$. See Remark~\ref{rmk:NS-precanonical-path}.
\end{enumerate}
Items~(a)--(d) have recorded pre-canonical derivations at
stated-unconditional level; canonical force requires
re-derivation in canonical papers.
Item~(e) records a conditional closure path pending
Theorem~\ref{thm:NS61}; the path has no canonical force.
\end{remark}

\begin{remark}[\status{R}]
\textup{(Current canonical re-derivation state.)}
The following results have now been re-derived in canonical papers:
\begin{itemize}
 \item NS Stage-0: $f_\infty > 0$ exists (Thm.~\ref{thm:NS-stage0},~[U]).
 \item NS Stage-1: $\varphi_\infty \geq 0$ exists (Thm.~\ref{thm:NS-stage1},~[U]).
 \item NS Stage-2a: Discrete balance identity (Thm.~\ref{thm:NS-stage2a},~[U]).
 \item NS Stage-2b: Nonlinear term identification (Thm.~\ref{thm:NS-stage2b},~[C]
 on KPO-3).
 \item Leray weak existence: global $L^2$ weak solution
 (Thm.~\ref{thm:NS-leray},~[U]).
\end{itemize}
NS Stage-3~(Thm.~\ref{thm:NS61}) remains~[O] in canonical papers.
\end{remark}

%% ---------------------------------------------------------------
\section{Boundary of the NS Branch Book}
%% ---------------------------------------------------------------

\textbf{What the NS branch book establishes.}
\begin{enumerate}
 \item The NS branch is a lawful descendant of $\mathbf{U}$,
 satisfying all five conditions of the Universal Branch
 Legitimacy Theorem.
 \item The NS obstruction quantity is well-defined and
 source-descended, not an external primitive.
 \item The bridge stack from source-descended obstruction
 control to the NS endpoint target is completely
 specified, with the missing step named precisely as
 NS.6.1.
 \item The proof program for NS.6.1 is structured and explicit,
 with the current frontier narrowed to UH.3 on the spike-free proof-program side and SB2 on the reserve-side structural side.
 \item Strong Boundary Safety is not the continuation theorem
 in disguise; it is strictly upstream of it.
\end{enumerate}

\textbf{What the NS branch book does not establish.}
\begin{enumerate}
 \item \emph{NS closure.} The NS Frontier Theorem is
 unconditionally proved; NS.6.1 is not.
 \item \emph{Proof of the full conditional package.} The branch still lacks discharge of the last proof-program micro-lemma UH.3 or explicit adoption of its spike-free-side dissipation replacement, and it still lacks unconditional reserve on the safe tail beyond SB2.
 \item \emph{Resolution of the NS Millennium Prize question.}
 That reduction is downstream of both UH.3/SF and SB2, together with the endpoint bridge NS.7.1.
 \item \emph{Any source-level claim.} No theorem in this
 branch book modifies or supplements the content of
 Books~I--VII.
\end{enumerate}

\begin{remark}[\status{R}]
For the branch's operative conditional theorem package, see Section~\ref{sec:NS-near-closure-modes}. For the final residual frontier after all reductions, see the current frontier note in Section~\ref{sec:NS-conditional-package} and the status statement in Section~12.
\end{remark}

%% ---------------------------------------------------------------
\section{Dependency Ledger}
%% ---------------------------------------------------------------

\renewcommand{\arraystretch}{1.4}
\noindent
\begin{tabular}{@{}p{3.8cm}p{0.8cm}p{3.2cm}p{5.2cm}@{}}
\hline
\textbf{Theorem} & \textbf{St.} & \textbf{Depends on}
 & \textbf{Exports to} \\
\hline
\ref{prop:NS-lawful}\ NS Lawful Descent
 & [U] & Bk.~IV, Bk.~V
 & \ref{prop:NS-witness}, \ref{thm:NS-frontier} \\

\ref{prop:NS-witness}\ NS Legitimacy Witness
 & [U] & Bk.~V
 & \ref{cor:NS-no-undeclared}, \ref{thm:NS-frontier} \\

\ref{thm:NS-obstruction-exists}\ Source-to-NS Obstr.
 & [U] & Bk.~III, Defs.~\ref{def:NS-active-weight}--\ref{def:NS-obstruction}
 & \ref{thm:NS61}, \ref{thm:NS-frontier} \\

\ref{prop:NS-criterion-required}\ Criterion Required
 & [U] & Bk.~VI
 & \ref{cor:NS-no-unnamed-criterion}, \ref{thm:NS61} \\

\ref{thm:NSP1}\ Window Coherence
 & [U] & Bk.~III, Def.~\ref{def:NS-contraction-profile}
 & \ref{thm:NSP1b} \\

\ref{thm:NSP1b}\ Successor Decomp.
 & [C] & \ref{thm:NSP1}, Bk.~III three-source
 & \ref{thm:NSP1c} \\

\ref{thm:NSL1a}\ Internal Defect Subord.
 & [C] & \ref{thm:NSP1b}, Phase-A floor
 & \ref{thm:NSL1} \\

\ref{thm:NSRen5}\ Late Align.-to-Dissip.
 & [C] & Active architecture
 & SB1 chain $\Rightarrow$ \ref{thm:NSL1b} \\

\ref{thm:NSSB2}\ Positive Reserve
 & [C] & Spike-free architecture
 & \ref{thm:NSG1plus} \\

\ref{thm:NSL1b}\ Spike-Res.\ Bdry Subord.
 & [C] & \ref{thm:NSRen5}, \ref{thm:NSSB2}
 & \ref{thm:NSL1} \\

\ref{thm:NSL1}\ Split-Defect Subord.
 & [C] & \ref{thm:NSL1a}, \ref{thm:NSL1b}
 & \ref{thm:NSP1c} \\

\ref{thm:NSG1plus}\ Strict Gap Est.
 & [C] & \ref{thm:NSL1}, transport ceiling
 & \ref{thm:NSP1c} \\

\ref{thm:NSP1c}\ Above-Threshold Contr.
 & [C] & \ref{thm:NSP1b}, \ref{thm:NSL1}
 & \ref{thm:NS61} \\

\ref{thm:NS61}\ Obstr.-to-Criterion Thm.
 & [C] & \ref{thm:NSP1}--\ref{thm:NSP1c},
 Def.~\ref{def:NS-active-criterion}
 & \ref{cor:NS61-consequence}, \ref{cor:NS-both-give-endpoint} \\

\ref{thm:NS71}\ Criterion-to-Endpoint
 & [C] & Def.~\ref{def:NS-active-criterion}, Bk.~VI
 & \ref{cor:NS-both-give-endpoint} \\

\ref{thm:NS-frontier}\ NS Frontier Thm.
 & [U] & \ref{prop:NS-lawful}, \ref{prop:NS-criterion-required},
 \ref{thm:NS61}, Bk.~VI, Bk.~VII
 & \ref{cor:NS-rephrasing-useless}, \ref{cor:NS-no-unconditional} \\

\ref{prop:NS-status}\ NS Status Prop.
 & [C] & \ref{prop:NS-lawful}, \ref{thm:NS-frontier}
 & Branch summary; external status statement \\
\hline
\end{tabular}
\renewcommand{\arraystretch}{1}

\medskip

\begin{scholium}[\status{R}]\label{sch:NS-final}
The NS branch is lawful, structurally serious, and precisely
organized. It is not blocked by a vague missing bridge. The
frontier items that once blocked it --- UH.3 on the proof-program
side and SB2 on the reserve-side structural side, each with a
precise mathematical statement, a clear logical role in the proof
graph, and an explicit anti-circularity argument --- have been
conditionally discharged, and the branch holds its domain-bounded
conditional endpoint claim at CERT-PROJ force
(Prop.~\ref{prop:NS-status}; MIG-052). What remains is a pair of
named residuals: the \emph{internal} obstruction/ledger
common-state bridge (Obl.~\ref{ob:NS-common-state-bridge}), through
which every transfer from the contracting ledger to the window
obstruction passes, and the \emph{external unconditional} frontier
--- unconditional global regularity (the Stage-3 Clay question) ---
which lies outside the declared conditional scope and is claimed
nowhere in this book.
That precision --- named items, named conditions, named residual ---
is the discipline Book~VII mandates, and the NS branch book
satisfies it.
\end{scholium}

\begin{remark}[\status{R}]\label{rem:NS-precanonical-completeness}
\textup{(Note on the pre-canonical record.)} This note concerns the
\emph{pre-canonical} v8.41 record and its relation to the canonical
body; the closing assessment of the canonical branch itself is
Scholium~\ref{sch:NS-final} above.
The pre-canonical program has recorded derivation sketches for
every structural sub-obligation except NS Stage-3 (the Clay Prize
question). These derivation sketches have no canonical force;
they identify which canonical proofs remain to be written. All
obligations within reach of the current NFC framework are discharged.
The frontier Theorem~\ref{thm:NS61} is not vague: it asks whether
the LR-weighted boundary-order condition
(Definition~\ref{def:NS-strong-ba}) holds, and the dream-suite
resolution path shows it does under $K_0=7$. That import is now
reflected in the canonical body: NS.6.1 is conditionally discharged
(Cor.~\ref{cor:NS61-discharged}), NS.7.1 is conditionally proved
(Thm.~\ref{thm:NS71}, Cor.~\ref{cor:NS-both-give-endpoint}), and the
branch holds its domain-bounded conditional endpoint claim at
CERT-PROJ force (Prop.~\ref{prop:NS-status}; MIG-052 --- the
transfer to the window obstruction passes through the open
common-state bridge, Obl.~\ref{ob:NS-common-state-bridge}).
Unconditional global regularity (NS Stage-3, the Clay Prize
question) and the internal common-state bridge are the named
residual frontiers.
\end{remark}

\end{document}
